% =====================================================================
%  Chương 2 — Các kỹ thuật học máy cơ bản (làn L2)
%  Mọi con số thực nghiệm lấy từ \SL{ch2_*}; mã trích từ src/TieuLuan 03/{tl3,scripts}/ch2_*.py
% =====================================================================
\providecommand{\bonhanxet}[3]{%
  \begin{phathien}{Quan sát -- Diễn giải -- Hàm ý cho mô hình}
  \textbf{Quan sát.} #1\par\smallskip
  \textbf{Diễn giải.} #2\par\smallskip
  \textbf{Hàm ý.} #3
  \end{phathien}}

\chapter{Các kỹ thuật học máy cơ bản (Machine Learning)}
\label{chap:ml}

Chương~1 cho thấy bước ngoặt của trí tuệ nhân tạo là chuyển từ \emph{viết tay tri thức} sang \emph{để máy học từ
dữ liệu}. Chương này trình bày nền tảng của bước chuyển đó: học máy (machine learning). Phần đầu (các mục
\ref{sec:ml-concepts}--\ref{sec:ml-scratch}) hệ thống hóa các khái niệm, thuật toán kinh điển và cách đánh giá mô hình,
mỗi thuật toán có công thức, hình minh họa và mã cài đặt. Phần sau là ba thực nghiệm đầy đủ theo cùng một quy trình:
dự đoán bệnh tiểu đường (phân loại nhị phân trên \SL{ch2_diabetes/survey/shape/0} hồ sơ), dự đoán giá nhà (hồi quy) và
phân loại nhận xét khách hàng (phân loại văn bản). Mỗi thực nghiệm đi qua đủ các bước: mô tả bài toán, giới thiệu dữ liệu,
khảo sát, làm sạch, biểu diễn, phân tích khám phá (EDA), chia tập, xây dựng mô hình, đánh giá, lựa chọn và triển khai.

Theo định nghĩa thường được trích dẫn của Mitchell~\cite{ch2:mitchell1997}: \emph{một chương trình máy tính được nói là
học từ kinh nghiệm $E$ đối với lớp nhiệm vụ $T$ và độ đo hiệu năng $P$, nếu hiệu năng của nó trên các nhiệm vụ thuộc
$T$, đo bằng $P$, được cải thiện nhờ kinh nghiệm $E$}. Với bài toán tiểu đường, $T$ là phân loại một hồ sơ bệnh nhân,
$E$ là các hồ sơ đã biết kết quả chẩn đoán và $P$ là F1 hay ROC-AUC trên những hồ sơ mô hình chưa từng thấy.

% =====================================================================
\section{Các khái niệm nền tảng}
\label{sec:ml-concepts}

\subsection{Dữ liệu, đặc trưng và nhãn}
Một tập dữ liệu có giám sát gồm $n$ cặp $\mathcal{D}=\{(\mathbf{x}_i, y_i)\}_{i=1}^{n}$. Mỗi \emph{mẫu} (sample,
observation) là một vector đặc trưng $\mathbf{x}_i\in\mathbb{R}^d$; mỗi thành phần $x_{ij}$ là một \emph{đặc trưng}
(feature, thuộc tính). Xếp các vector thành hàng ta được \emph{ma trận thiết kế} $X\in\mathbb{R}^{n\times d}$. Biến
cần dự đoán $y_i$ là \emph{nhãn} (label, target): rời rạc với bài toán phân loại (classification), liên tục với bài toán
hồi quy (regression).

Đặc trưng có hai dạng chính. \emph{Đặc trưng số} (numerical) như tuổi, BMI, diện tích có phép so sánh và phép trừ có
nghĩa. \emph{Đặc trưng phân loại} (categorical) như giới tính, thành phố chỉ là tên nhóm; muốn đưa vào mô hình phải mã
hóa thành số (mục~\ref{sec:ml-feature}). Văn bản, ảnh và chuỗi thời gian là dữ liệu phi cấu trúc, cần bước biểu diễn
riêng: TF-IDF cho văn bản ở chương này, tích chập cho ảnh ở Chương~3, mạng hồi quy cho chuỗi ở Chương~4.

\begin{dinhnghia}{Học máy như bài toán xấp xỉ hàm}
Giả sử tồn tại hàm chưa biết $f:\mathcal{X}\to\mathcal{Y}$ sinh ra nhãn (cộng nhiễu). Học máy chọn trong một
\emph{lớp giả thuyết} $\mathcal{H}$ (ví dụ mọi hàm tuyến tính, mọi cây độ sâu $\le 8$) một hàm $h\in\mathcal{H}$ sao cho
$h(\mathbf{x})\approx y$ \emph{trên dữ liệu mới}, không chỉ trên dữ liệu đã thấy.
\end{dinhnghia}

\subsection{Ba kiểu học chính}
Hình~\ref{fig:ch2-types} tóm tắt ba kiểu học. \textbf{Học có giám sát} (supervised learning) dùng dữ liệu có nhãn để
học ánh xạ $\mathbf{x}\mapsto y$; đây là kiểu của cả ba thực nghiệm trong chương. \textbf{Học không giám sát}
(unsupervised learning) chỉ có $\mathbf{x}$ và tìm cấu trúc ẩn như cụm (K-means) hay các hướng biến thiên chính (PCA),
được minh họa ở mục~\ref{sec:ml-unsup}. \textbf{Học tăng cường} (reinforcement learning) học một chính sách hành động
từ phần thưởng trả về sau mỗi tương tác với môi trường; đây là nền tảng của AlphaGo được nhắc ở Chương~1.

\hinhtikz{ch2_learning_types}{Ba kiểu học máy: có giám sát (đường biên tách hai lớp có nhãn), không giám sát (tự tìm
ba cụm) và tăng cường (vòng lặp tác tử -- môi trường với hành động $a_t$, trạng thái $s_{t+1}$, phần thưởng $r_{t+1}$).}{ch2-types}

\subsection{Hàm mất mát và huấn luyện}
\label{sec:ml-loss}
Để đo mức sai của $h$ trên một mẫu, ta dùng \emph{hàm mất mát} $\ell(h(\mathbf{x}), y)$. Huấn luyện là cực tiểu
\emph{rủi ro thực nghiệm} (empirical risk), thường kèm một số hạng chính quy hóa $\Omega$ phạt mô hình phức tạp:
\begin{equation}
  \hat{\mathbf{w}}=\arg\min_{\mathbf{w}}\; J(\mathbf{w})=\frac{1}{n}\sum_{i=1}^{n}\ell\big(h_{\mathbf{w}}(\mathbf{x}_i),y_i\big)+\lambda\,\Omega(\mathbf{w}).
  \label{eq:erm}
\end{equation}
Hai hàm mất mát dùng nhiều nhất là sai số bình phương $\ell=(\hat{y}-y)^2$ cho hồi quy và cross-entropy
$\ell=-[y\log p+(1-y)\log(1-p)]$ cho phân loại. Khi không có nghiệm đóng, ta dùng \textbf{Gradient Descent}: lặp lại
\begin{equation}
  \mathbf{w}^{(t+1)}=\mathbf{w}^{(t)}-\eta\,\nabla_{\mathbf{w}}J\big(\mathbf{w}^{(t)}\big),
\end{equation}
với $\eta$ là tốc độ học (learning rate). Phiên bản \emph{ngẫu nhiên} (SGD) và \emph{mini-batch} ước lượng gradient
trên một phần nhỏ dữ liệu mỗi bước, rẻ hơn nhiều với $n$ lớn; Adam~\cite{ch2:kingma2015} điều chỉnh $\eta$ riêng cho
từng tham số theo mô-men bậc một và bậc hai của gradient. Hình~\ref{fig:ch2-gd} cho thấy vai trò của $\eta$ trên một
hàm mất mát bậc hai lệch trục: $\eta$ quá nhỏ hội tụ chậm, $\eta$ vừa phải đi thẳng vào cực tiểu, $\eta$ lớn gây dao
động qua lại theo trục dốc.

\hinhpdf{ch2_gd_lr}{1}{Đường đi của Gradient Descent trên mặt mất mát $J(\mathbf{w})=\tfrac12\mathbf{w}^\top A\mathbf{w}$
(ngôi sao là cực tiểu) với ba tốc độ học. Mặt mất mát dốc theo $w_1$ và thoải theo $w_2$ nên $\eta$ lớn dao động theo $w_1$.}{ch2-gd}

\subsection{Chia dữ liệu: huấn luyện, kiểm định, kiểm thử}
Mục tiêu của học máy là \emph{tổng quát hóa}. Vì vậy dữ liệu được chia thành ba phần không giao nhau
(Hình~\ref{fig:ch2-split}a): \textbf{Train} để học tham số, \textbf{Validation} để so sánh mô hình, chọn siêu tham số
và ngưỡng quyết định, \textbf{Test} để báo cáo hiệu năng cuối cùng và chỉ dùng một lần. Với bài toán phân loại mất cân
bằng, việc chia được \emph{phân tầng} (stratified) để tỷ lệ các lớp ở ba tập như nhau. Khi dữ liệu ít, \emph{kiểm định
chéo $k$-fold}~\cite{ch2:kohavi1995} (Hình~\ref{fig:ch2-split}b) cho ước lượng ổn định hơn: dữ liệu chia $k$ phần, mỗi
lượt dùng một phần làm kiểm định và $k-1$ phần còn lại để huấn luyện, rồi lấy trung bình và độ lệch chuẩn của $k$ điểm số.

\hinhtikz{ch2_split}{(a) Chia hold-out 70/15/15 dùng trong cả ba thực nghiệm; (b) kiểm định chéo 5-fold trên tập
Train dùng để ước lượng độ ổn định của mô hình.}{ch2-split}

\begin{luuy}{Rò rỉ dữ liệu (data leakage)}
Mọi phép biến đổi học tham số từ dữ liệu (trung bình và độ lệch chuẩn của chuẩn hóa, từ vựng TF-IDF, ngưỡng quyết
định) phải được \emph{fit trên Train}, rồi chỉ \emph{áp dụng} cho Validation và Test. Trong mã của chương, điều này được
bảo đảm bằng cách gói tiền xử lý và mô hình vào cùng một \texttt{Pipeline} của scikit-learn~\cite{ch2:pedregosa2011}.
\end{luuy}

% =====================================================================
\section{Quy trình học máy điển hình}
Hình~\ref{fig:ch2-pipeline} là quy trình tám bước mà cả ba thực nghiệm tuân theo, cũng là cấu trúc đề mục của các
mục~\ref{sec:ml-diab}--\ref{sec:ml-cmt}. Điểm đáng chú ý là vòng phản hồi: kết quả đánh giá trên Validation thường đưa
ta quay lại thiết kế đặc trưng hoặc siêu tham số. Ranh giới quan trọng là tập Test không được tham gia vòng lặp này.

\hinhtikz{ch2_pipeline}{Quy trình học máy có giám sát tám bước, từ thu thập dữ liệu tới triển khai; vòng đứt nét là
phản hồi từ đánh giá trên Validation về đặc trưng và siêu tham số.}{ch2-pipeline}

\begin{enumerate}
  \item \textbf{Thu thập dữ liệu} và hiểu ngữ cảnh: nguồn, đơn vị, ý nghĩa từng cột.
  \item \textbf{Khảo sát và làm sạch}: giá trị thiếu, bản ghi trùng, giá trị vô lý, ngoại lai.
  \item \textbf{EDA và biểu diễn}: phân bố, tương quan, quan hệ với nhãn; mã hóa và chuẩn hóa.
  \item \textbf{Chia tập} Train/Validation/Test, phân tầng nếu là phân loại.
  \item \textbf{Huấn luyện} nhiều họ mô hình, luôn kèm một \emph{baseline} ngây thơ để biết mức ``không học gì''.
  \item \textbf{Đánh giá và tinh chỉnh} bằng các chỉ số phù hợp bản chất bài toán.
  \item \textbf{Chọn mô hình} dựa trên Validation, rồi đánh giá một lần trên Test.
  \item \textbf{Triển khai}: đóng gói tiền xử lý + mô hình thành dịch vụ dự đoán.
\end{enumerate}

% =====================================================================
\section{Hồi quy tuyến tính (Linear Regression)}
\label{sec:ml-linreg}

\subsection{Bình phương tối thiểu và hàm mất mát}
Hồi quy tuyến tính giả định nhãn là tổ hợp tuyến tính của đặc trưng cộng nhiễu:
$y=w_0+w_1x_1+\dots+w_dx_d+\varepsilon=\tilde{\mathbf{x}}^\top\mathbf{w}+\varepsilon$, trong đó
$\tilde{\mathbf{x}}=(1,x_1,\dots,x_d)$ đã thêm thành phần 1 cho hệ số chặn. Phương pháp bình phương tối thiểu (OLS)
chọn $\mathbf{w}$ cực tiểu trung bình bình phương phần dư (Hình~\ref{fig:ch2-linreg}):
\begin{equation}
  J(\mathbf{w})=\frac{1}{n}\lVert X\mathbf{w}-\mathbf{y}\rVert_2^2,\qquad
  \nabla J=\frac{2}{n}X^\top(X\mathbf{w}-\mathbf{y})=0\;\Rightarrow\;\hat{\mathbf{w}}=(X^\top X)^{-1}X^\top\mathbf{y}.
  \label{eq:ols}
\end{equation}
Công thức (\ref{eq:ols}) là \emph{phương trình chuẩn} (normal equation). Nó cho nghiệm chính xác trong một bước nhưng
chi phí $O(nd^2+d^3)$; khi $d$ lớn hoặc $X^\top X$ gần suy biến, ta dùng Gradient Descent hoặc phân rã QR/SVD.

\hinhtikz{ch2_linreg}{Hồi quy tuyến tính một biến: đường $\hat y=w_0+w_1x$ cực tiểu tổng bình phương phần dư (đoạn
đứt nét màu cam); bên phải là ba cách tìm $\mathbf{w}$.}{ch2-linreg}

\textbf{Ví dụ số nhỏ.} Với ba điểm $(1,2),(2,3),(3,5)$: $\bar{x}=2$, $\bar{y}=10/3$; hệ số góc
$w_1=\sum(x_i-\bar x)(y_i-\bar y)/\sum(x_i-\bar x)^2=(-1)(-4/3)+0+(1)(5/3)$ chia cho $2$, tức $w_1=1{,}5$, và
$w_0=\bar y-w_1\bar x=1/3$. Đường thẳng $\hat{y}=0{,}33+1{,}5x$ có phần dư $0{,}17;\,-0{,}33;\,0{,}17$, tổng bằng 0 đúng
như tính chất của OLS khi có hệ số chặn.

Đoạn mã~\ref{code:ch2_scratch__linreg_normal_equation} và~\ref{code:ch2_scratch__linreg_gd} cài đặt cả hai cách
bằng NumPy; mục~\ref{sec:ml-scratch} đối chiếu chúng với scikit-learn trên dữ liệu giá nhà.

\maNguon{ch2_scratch__linreg_normal_equation}{Hồi quy tuyến tính bằng phương trình chuẩn (NumPy)}
\maNguon{ch2_scratch__linreg_gd}{Hồi quy tuyến tính bằng Gradient Descent trên MSE (NumPy)}

\subsection{Mở rộng: đa thức, Ridge và Lasso}
\textbf{Hồi quy đa thức} vẫn là hồi quy tuyến tính theo tham số, chỉ thay $\mathbf{x}$ bằng các đặc trưng mở rộng
$\phi(\mathbf{x})=(x_1,x_2,x_1^2,x_1x_2,x_2^2,\dots)$, nhờ đó học được quan hệ cong. Bậc càng cao, mô hình càng linh hoạt
và càng dễ quá khớp (mục~\ref{sec:ml-overfit}). \textbf{Ridge}~\cite{ch2:hoerl1970} thêm phạt L2, \textbf{Lasso}%
~\cite{ch2:tibshirani1996} thêm phạt L1:
\begin{equation}
  J_{\text{Ridge}}=\tfrac{1}{n}\lVert X\mathbf{w}-\mathbf{y}\rVert^2+\lambda\lVert\mathbf{w}\rVert_2^2,\qquad
  J_{\text{Lasso}}=\tfrac{1}{2n}\lVert X\mathbf{w}-\mathbf{y}\rVert^2+\lambda\lVert\mathbf{w}\rVert_1 .
\end{equation}
Ridge vẫn có nghiệm đóng $\hat{\mathbf{w}}=(X^\top X+n\lambda I)^{-1}X^\top\mathbf{y}$, luôn khả nghịch nên ổn định khi
các đặc trưng tương quan. Lasso không có nghiệm đóng nhưng có tính chất quý: đưa một số hệ số về \emph{đúng bằng 0},
tức tự chọn đặc trưng. Hình~\ref{fig:ch2-l1l2} giải thích bằng hình học: nghiệm là điểm đường đồng mức của hàm mất mát
chạm miền ràng buộc lần đầu; miền L1 là hình thoi có đỉnh nằm trên trục nên điểm chạm thường là đỉnh, nơi một tọa độ bằng 0.

\hinhtikz{ch2_l1l2}{Hình học của chính quy hóa: đường đồng mức (cam) của MSE quanh nghiệm OLS chạm miền ràng buộc. Hình
tròn L2 cho nghiệm co nhỏ đều; hình thoi L1 cho nghiệm nằm ở đỉnh, khiến $w_1=0$.}{ch2-l1l2}

% =====================================================================
\section{Perceptron, hồi quy logistic và Softmax}
\label{sec:ml-logreg}

\subsection{Perceptron Learning Algorithm (PLA)}
Perceptron của Rosenblatt~\cite{ch2:rosenblatt1958} là bộ phân loại tuyến tính sớm nhất có thuật toán học:
$\hat{y}=\operatorname{sign}(\mathbf{w}^\top\tilde{\mathbf{x}})$ với $y\in\{-1,+1\}$. Mỗi khi gặp điểm bị phân loại sai
($y_i\mathbf{w}^\top\tilde{\mathbf{x}}_i\le 0$), PLA cập nhật $\mathbf{w}\leftarrow\mathbf{w}+y_i\tilde{\mathbf{x}}_i$,
tức xoay siêu phẳng về phía điểm đó. Định lý hội tụ Novikoff bảo đảm PLA dừng sau hữu hạn bước \emph{nếu} dữ liệu tách
được tuyến tính; ngược lại nó dao động mãi, đó cũng là giới hạn XOR đã gây ra mùa đông AI thứ nhất (Chương~1).
Trên hai cụm tách được (Hình~\ref{fig:ch2-pla}, trái), cài đặt ở Đoạn mã~\ref{code:ch2_scratch__pla} hội tụ sau
\SL{ch2_scratch/pla/epochs} epoch.

\maNguon{ch2_scratch__pla}{Perceptron Learning Algorithm cài đặt bằng NumPy}

\subsection{Hồi quy logistic và hàm mất mát}
Thay vì nhãn cứng, \textbf{hồi quy logistic}~\cite{ch2:cox1958} ước lượng xác suất lớp dương bằng hàm sigmoid
(Hình~\ref{fig:ch2-logistic}):
\begin{equation}
  p=P(y=1\mid\mathbf{x})=\sigma(z)=\frac{1}{1+e^{-z}},\qquad z=\mathbf{w}^\top\mathbf{x}+b .
\end{equation}
Logarit tỷ số odds tuyến tính theo $\mathbf{x}$: $\log\frac{p}{1-p}=\mathbf{w}^\top\mathbf{x}+b$, nên mỗi hệ số $w_j$
có diễn giải trực tiếp: tăng $x_j$ một đơn vị làm odds nhân với $e^{w_j}$. Tham số học bằng cực đại hợp lý, tương đương
cực tiểu cross-entropy nhị phân (BCE). Gradient có dạng gọn đẹp
\begin{equation}
  \nabla_{\mathbf{w}}\mathcal{L}=\frac{1}{n}X^\top\big(\sigma(X\mathbf{w})-\mathbf{y}\big),
\end{equation}
giống hệt hồi quy tuyến tính với $\hat{\mathbf{y}}$ thay bằng $\sigma(X\mathbf{w})$. BCE lồi theo $\mathbf{w}$ nên GD
hội tụ về cực tiểu toàn cục. Hình~\ref{fig:ch2-sigmoid} cho thấy BCE phạt rất nặng các dự đoán tự tin nhưng sai.

\hinhtikz{ch2_logistic}{Hồi quy logistic như một nơ-ron: tổng có trọng số $z$, hàm sigmoid cho xác suất $p$, so với
ngưỡng $\tau$ để ra nhãn.}{ch2-logistic}
\hinhpdf{ch2_sigmoid_logloss}{1}{Trái: hàm sigmoid ánh xạ $z\in\mathbb{R}$ vào $(0,1)$. Phải: cross-entropy nhị phân
tiến tới vô cùng khi mô hình tự tin sai.}{ch2-sigmoid}

\maNguon{ch2_scratch__logreg_gd}{Hồi quy logistic bằng mini-batch Gradient Descent trên BCE + L2 (NumPy)}

\subsection{Ngưỡng quyết định}
Nhãn dự đoán là $\hat{y}=\mathbf{1}[p\ge\tau]$. Ngưỡng mặc định $\tau=0{,}5$ chỉ tối ưu cho accuracy khi hai loại lỗi
có chi phí như nhau và các lớp cân bằng. Với dữ liệu y tế mất cân bằng, bỏ sót bệnh (FN) đắt hơn báo động nhầm (FP),
nên $\tau$ được chọn trên tập Validation để cực đại F1 hoặc đạt Recall mong muốn (Đoạn mã~\ref{code:ch2_metrics__best_threshold}).

\subsection{Mở rộng nhiều lớp: Softmax Regression}
Với $K$ lớp, mỗi lớp có một vector trọng số $\mathbf{w}_k$ và xác suất được chuẩn hóa bằng hàm softmax:
\begin{equation}
  \begin{gathered}
  P(y=k\mid\mathbf{x})=\frac{\exp(\mathbf{w}_k^\top\mathbf{x})}{\sum_{j=1}^{K}\exp(\mathbf{w}_j^\top\mathbf{x})},\\
  \mathcal{L}=-\frac{1}{n}\sum_{i}\sum_{k}\mathbf{1}[y_i=k]\log P(y_i=k\mid\mathbf{x}_i).
  \end{gathered}
\end{equation}
Gradient theo ma trận $W$ là $X^\top(P-Y)/n$, với $Y$ là ma trận one-hot của nhãn. Hình~\ref{fig:ch2-pla} (phải) là
biên quyết định Softmax 3 lớp do Đoạn mã~\ref{code:ch2_scratch__softmax_gd} học được, đạt accuracy
\SL{ch2_scratch/softmax/acc}. Softmax cũng là lớp ra của mọi mạng phân loại nhiều lớp ở Chương~3. Trong TN3, Softmax
Regression được dùng để dự đoán số sao (1--5) của một nhận xét.

\hinhpdf{ch2_pla_softmax}{1}{Trái: siêu phẳng do PLA tìm được trên hai cụm tách được tuyến tính. Phải: ba vùng quyết định
của Softmax Regression; các biên giữa hai vùng là đường thẳng.}{ch2-pla}
\maNguon{ch2_scratch__softmax_gd}{Softmax Regression $K$ lớp bằng Gradient Descent (NumPy)}

% =====================================================================
\section{Đánh giá mô hình}
\label{sec:ml-eval}

\subsection{Ma trận nhầm lẫn và các chỉ số phân loại}
Với phân loại nhị phân, mọi dự đoán rơi vào một trong bốn ô của \emph{ma trận nhầm lẫn} (Hình~\ref{fig:ch2-confusion}).
Accuracy là tỷ lệ đúng tổng thể nhưng \emph{gây hiểu lầm khi mất cân bằng}: trên tập tiểu đường, mô hình luôn đoán
``không bệnh'' đạt accuracy \SL{ch2_diabetes/val/baseline/accPct}\% mà không phát hiện được ca nào. Precision đo độ tin
cậy của các báo động, Recall (độ nhạy) đo tỷ lệ ca bệnh được phát hiện, F1 là trung bình điều hòa của hai đại lượng.
Với nhiều lớp, \emph{macro-F1} lấy trung bình F1 của từng lớp, coi các lớp quan trọng như nhau.

\hinhtikz{ch2_confusion}{Ma trận nhầm lẫn và bốn chỉ số phân loại suy ra từ nó.}{ch2-confusion}

\subsection{Đường ROC và Precision--Recall}
Khi quét ngưỡng $\tau$ từ 1 xuống 0, ta được đường \textbf{ROC} (TPR theo FPR); diện tích dưới đường (ROC-AUC) bằng xác
suất mô hình xếp một mẫu dương ngẫu nhiên cao hơn một mẫu âm ngẫu nhiên~\cite{ch2:fawcett2006}. ROC-AUC không phụ thuộc
ngưỡng và không phụ thuộc tỷ lệ lớp, nên lại có thể ``lạc quan'' khi lớp dương hiếm. Đường \textbf{Precision--Recall}
và diện tích của nó (PR-AUC, hay Average Precision) nhạy hơn với hiệu năng trên lớp hiếm~\cite{ch2:saito2015}: mức
ngẫu nhiên của PR-AUC chính là tỷ lệ lớp dương, chứ không phải 0,5.

\subsection{Chỉ số cho hồi quy}
Với $e_i=y_i-\hat{y}_i$:
\begin{equation}
  \begin{gathered}
  \mathrm{MAE}=\frac{1}{n}\sum|e_i|,\qquad \mathrm{RMSE}=\sqrt{\frac{1}{n}\sum e_i^2},\\
  R^2=1-\frac{\sum e_i^2}{\sum(y_i-\bar{y})^2},\qquad \mathrm{MAPE}=\frac{1}{n}\sum\left|\frac{e_i}{y_i}\right| .
  \end{gathered}
\end{equation}
MAE dễ diễn giải (cùng đơn vị với $y$) và ít nhạy với ngoại lai; RMSE phạt nặng sai số lớn; $R^2$ là tỷ lệ phương sai
của $y$ được giải thích, bằng 0 với mô hình luôn dự đoán trung bình; MAPE là sai số tương đối, thuận tiện khi so sánh
giữa các thang giá.

\maNguon{ch2_metrics__clf_metrics}{Tính các chỉ số phân loại ở một ngưỡng cho trước}
\maNguon{ch2_metrics__best_threshold}{Chọn ngưỡng quyết định theo F1 trên tập Validation}
\maNguon{ch2_metrics__reg_metrics}{Bốn chỉ số hồi quy MAE, RMSE, $R^2$, MAPE}

% =====================================================================
\section{Quá khớp, chính quy hóa và thiên lệch -- phương sai}
\label{sec:ml-overfit}

\subsection{Quá khớp và chưa khớp}
Mô hình \emph{chưa khớp} (underfitting) quá đơn giản để nắm được quy luật: sai số cao trên cả Train lẫn Test. Mô hình
\emph{quá khớp} (overfitting) học thuộc cả nhiễu: sai số Train rất thấp nhưng Test cao. Hình~\ref{fig:ch2-poly} minh họa
bằng đa thức bậc 1, 4 và 15 khớp 30 điểm sinh từ hàm sin cộng nhiễu.

\hinhpdf{ch2_poly_fit}{1}{Đa thức bậc 1 (chưa khớp), bậc 4 (khớp vừa) và bậc 15 (quá khớp, dao động giữa các điểm) trên cùng
30 điểm dữ liệu; đường đứt nét là hàm thật.}{ch2-poly}

\subsection{Phân rã thiên lệch -- phương sai}
Với sai số bình phương, sai số kỳ vọng tại một điểm phân rã thành ba thành phần~\cite{ch2:geman1992}:
\begin{equation}
  \mathbb{E}\big[(y-\hat{f}(\mathbf{x}))^2\big]=\underbrace{\big(\mathbb{E}[\hat f(\mathbf{x})]-f(\mathbf{x})\big)^2}_{\text{thiên lệch}^2}
  +\underbrace{\mathbb{E}\big[(\hat f(\mathbf{x})-\mathbb{E}[\hat f(\mathbf{x})])^2\big]}_{\text{phương sai}}+\underbrace{\sigma^2}_{\text{nhiễu}} .
\end{equation}
Mô hình đơn giản có thiên lệch cao, phương sai thấp; mô hình phức tạp ngược lại (Hình~\ref{fig:ch2-dart}). Sai số Test
có dạng chữ U theo độ phức tạp (Hình~\ref{fig:ch2-bv}); chọn mô hình là tìm đáy của chữ U đó bằng Validation.

\hinhtikz{ch2_bias_dart}{Bốn tổ hợp thiên lệch và phương sai: tâm bia là giá trị thật, mỗi chấm là dự đoán của một mô
hình huấn luyện trên một tập dữ liệu khác nhau.}{ch2-dart}
\hinhpdf{ch2_bias_variance}{1}{Sai số Train giảm đơn điệu theo bậc đa thức, còn sai số Test có dạng chữ U: vùng trái thiên
lệch cao, vùng phải phương sai cao.}{ch2-bv}

\subsection{Các cách giảm quá khớp}
\begin{itemize}
  \item \textbf{Chính quy hóa} L2/L1 (Ridge, Lasso, \texttt{weight\_decay} trong mạng nơ-ron); hệ số $C=1/\lambda$ của
        logistic và SVM.
  \item \textbf{Giới hạn độ phức tạp}: độ sâu cây, số lá tối thiểu, số láng giềng $k$ lớn hơn.
  \item \textbf{Tổ hợp mô hình}: Random Forest lấy trung bình nhiều cây để giảm phương sai.
  \item \textbf{Dừng sớm} (early stopping) theo mất mát Validation và \textbf{Dropout}~\cite{ch2:srivastava2014} với mạng nơ-ron.
  \item \textbf{Thêm dữ liệu}: đường học (learning curve) cho biết thêm dữ liệu còn giúp được hay không.
\end{itemize}

% =====================================================================
\section{Cây quyết định (Decision Trees)}
\label{sec:ml-tree}
Cây quyết định~\cite{ch2:quinlan1986,ch2:breiman1984} chia không gian đặc trưng bằng một chuỗi câu hỏi dạng
``$x_j\le t$?''. Tại mỗi nút, thuật toán CART thử mọi đặc trưng và ngưỡng, chọn phép tách làm giảm độ vẩn đục (impurity)
nhiều nhất, rồi đệ quy cho đến khi gặp điều kiện dừng (độ sâu tối đa, số mẫu tối thiểu ở lá). Với tỷ lệ lớp $p_k$ trong
một nút:
\begin{equation}
  \begin{gathered}
  \text{Gini}=1-\sum_k p_k^2,\qquad \text{Entropy}=-\sum_k p_k\log_2 p_k,\\
  \Delta=I(\text{cha})-\sum_{c}\frac{n_c}{n}I(\text{con}_c) .
  \end{gathered}
\end{equation}
Hai độ đo có hình dạng gần như nhau (Hình~\ref{fig:ch2-impurity}), cực đại khi hai lớp cân bằng và bằng 0 khi nút thuần.
Hình~\ref{fig:ch2-tree} là cây độ sâu 2 học trên dữ liệu tiểu đường thật (Đoạn mã~\ref{code:ch2_tree_demo__describe_tree}):
câu hỏi đầu tiên mà thuật toán tự tìm ra là HbA1c, đúng là xét nghiệm chẩn đoán lâm sàng~\cite{ch2:ada2023}.

\hinhpdf{ch2_impurity}{1}{Entropy, chỉ số Gini và lỗi phân loại theo tỷ lệ lớp dương trong một nút.}{ch2-impurity}
\hinhtikz{ch2_tree}{Cây quyết định độ sâu 2 huấn luyện trên \SL{ch2_diabetes/clean/n_rows} hồ sơ tiểu đường; mọi ngưỡng,
Gini và tỷ lệ mẫu đọc từ cây đã khớp.}{ch2-tree}
\maNguon{ch2_tree_demo__describe_tree}{Duyệt cấu trúc cây scikit-learn để lấy ngưỡng, Gini và xác suất lá}

Ưu điểm của cây là dễ diễn giải, không cần chuẩn hóa, xử lý được quan hệ phi tuyến và tương tác. Nhược điểm là phương
sai cao: một thay đổi nhỏ trong dữ liệu có thể đổi cả cây, và cây sâu dễ quá khớp. Đó là lý do ra đời các phương pháp tổ hợp.

% =====================================================================
\section{Phương pháp tổ hợp (Ensembles)}
\label{sec:ml-ensemble}
\textbf{Bagging}~\cite{ch2:breiman1996} huấn luyện $B$ mô hình trên $B$ mẫu bootstrap (lấy mẫu có hoàn lại) rồi lấy
trung bình. Nếu các mô hình có phương sai $\sigma^2$ và tương quan $\rho$, phương sai của trung bình là
$\rho\sigma^2+\frac{1-\rho}{B}\sigma^2$. \textbf{Random Forest}~\cite{ch2:breiman2001} giảm thêm $\rho$ bằng cách chỉ xét
một tập con ngẫu nhiên đặc trưng ở mỗi lần tách. \textbf{Boosting} xây cây \emph{tuần tự}: cây sau học phần sai của
tổng các cây trước. Gradient Boosting~\cite{ch2:friedman2001} khớp cây thứ $m$ với gradient âm của hàm mất mát
$g_i^{(m)}=-\partial\ell(y_i,F_{m-1}(\mathbf{x}_i))/\partial F$ rồi cập nhật $F_m=F_{m-1}+\eta f_m$.
\textbf{XGBoost}~\cite{ch2:chen2016} dùng khai triển Taylor bậc hai và phạt độ phức tạp của cây; \textbf{LightGBM}%
~\cite{ch2:ke2017} tăng tốc bằng histogram đặc trưng và mọc cây theo lá. Hình~\ref{fig:ch2-ensemble} so sánh hai cơ chế.

\hinhtikz{ch2_ensemble}{Bagging (trái): các cây độc lập trên mẫu bootstrap, kết quả bỏ phiếu. Boosting (phải): mỗi cây
học phần dư/gradient của tổng các cây trước, kết quả là tổng có hệ số học $\eta$.}{ch2-ensemble}

% =====================================================================
\section{Các thuật toán kinh điển khác}
\label{sec:ml-other}

\subsection{K láng giềng gần nhất (KNN)}
KNN~\cite{ch2:cover1967} không có bước huấn luyện: lưu toàn bộ dữ liệu, và khi dự đoán thì tìm $k$ điểm gần nhất rồi bỏ
phiếu (phân loại) hoặc lấy trung bình (hồi quy), có thể đánh trọng số $1/d$ (Hình~\ref{fig:ch2-knn}). Vì dựa trên khoảng
cách nên KNN \emph{bắt buộc phải chuẩn hóa} đặc trưng; nếu không, glucose (đơn vị hàng trăm) sẽ lấn át HbA1c (đơn vị vài
phần trăm). $k$ nhỏ cho biên quyết định răng cưa (phương sai cao), $k$ lớn làm biên mượt dần (Hình~\ref{fig:ch2-knnk}).
Chi phí dự đoán $O(nd)$ cho mỗi truy vấn là nhược điểm lớn với dữ liệu nhiều mẫu.

\hinhtikz{ch2_knn}{KNN: với $k=3$ điểm truy vấn (ngôi sao) được gán lớp đỏ, với $k=7$ lại thành lớp xanh: lựa chọn $k$
thay đổi dự đoán.}{ch2-knn}
\hinhpdf{ch2_knn_k}{1}{Biên quyết định của KNN trên dữ liệu ``mặt trăng'' với $k=1$, 15 và 75.}{ch2-knnk}

\subsection{Naive Bayes}
Naive Bayes áp dụng định lý Bayes với giả định ``ngây thơ'' rằng các đặc trưng độc lập có điều kiện khi biết lớp:
\begin{equation}
  P(y=k\mid\mathbf{x})\propto P(y=k)\prod_{j=1}^{d}P(x_j\mid y=k) .
\end{equation}
Gaussian NB mô hình $P(x_j\mid y)$ bằng phân phối chuẩn; Multinomial NB~\cite{ch2:mccallum1998} mô hình tần suất từ,
rất hợp với văn bản; Complement NB~\cite{ch2:rennie2003} ước lượng trên phần bù của lớp nên bền hơn khi mất cân bằng.
Mô hình huấn luyện cực nhanh nhưng xác suất đầu ra thường bị đẩy về 0 hoặc 1 do giả định độc lập bị vi phạm.

\subsection{Máy vector hỗ trợ (SVM) và SVR}
SVM~\cite{ch2:cortes1995} tìm siêu phẳng $\mathbf{w}^\top\mathbf{x}+b=0$ có \emph{lề} lớn nhất tới các điểm gần nhất
của hai lớp (các \emph{vector hỗ trợ}). Lề hình học là $2/\lVert\mathbf{w}\rVert$, dẫn đến bài toán tối ưu lồi:
\begin{equation}
  \min_{\mathbf{w},b,\boldsymbol\xi}\;\tfrac12\lVert\mathbf{w}\rVert^2+C\sum_i\xi_i\quad
  \text{s.t.}\quad y_i(\mathbf{w}^\top\mathbf{x}_i+b)\ge 1-\xi_i,\;\xi_i\ge0 .
\end{equation}
Ta có ba cấp độ: \emph{lề cứng} khi dữ liệu tách được ($\xi_i=0$), \emph{lề mềm} cho phép vi phạm với chi phí $C$
(Hình~\ref{fig:ch2-svm}), và \emph{kernel} cho dữ liệu không tách tuyến tính. Ở dạng đối ngẫu, dữ liệu chỉ xuất hiện qua
tích vô hướng $\mathbf{x}_i^\top\mathbf{x}_j$, có thể thay bằng hàm nhân $K(\mathbf{x}_i,\mathbf{x}_j)=\phi(\mathbf{x}_i)^\top
\phi(\mathbf{x}_j)$ như RBF $K=\exp(-\gamma\lVert\mathbf{x}_i-\mathbf{x}_j\rVert^2)$. Đây là ``kernel trick'': làm việc
ngầm trong không gian đặc trưng nhiều chiều mà không cần tính $\phi$ tường minh (Hình~\ref{fig:ch2-kernel}).
\textbf{SVR}~\cite{ch2:drucker1997} dùng cùng ý tưởng cho hồi quy: chỉ phạt những điểm nằm ngoài một ống rộng
$\pm\varepsilon$ quanh hàm dự đoán (hàm mất mát $\varepsilon$-insensitive).

\hinhpdf{ch2_svm_margin}{1}{SVM tuyến tính: $C$ lớn (trái) cho lề hẹp, ít vector hỗ trợ; $C$ nhỏ (phải) chấp nhận vi phạm
để có lề rộng, nhiều vector hỗ trợ (vòng cam).}{ch2-svm}
\hinhtikz{ch2_svm_kernel}{Kernel trick: hai vòng tròn đồng tâm không tách được trong $\mathbb{R}^2$ trở nên tách được
bằng một mặt phẳng sau phép nâng $\phi(\mathbf{x})=(x_1,x_2,x_1^2+x_2^2)$.}{ch2-kernel}

\subsection{So sánh trực quan các biên quyết định}
Hình~\ref{fig:ch2-boundaries} đặt chín mô hình lên cùng một bài toán hai chiều phi tuyến. Logistic Regression và SVM
tuyến tính chỉ vẽ được đường thẳng; Naive Bayes cho biên cong bậc hai; cây quyết định cho các bậc thang song song trục;
Random Forest và XGBoost mượt hơn nhờ trung bình nhiều cây; KNN, SVM RBF và logistic với đặc trưng đa thức bám theo hình
dạng thật của dữ liệu. Mỗi họ mô hình mang một \emph{thiên kiến quy nạp} (inductive bias) khác nhau, và không có họ nào
tốt nhất cho mọi bài toán.

\hinhpdf{ch2_boundaries}{1}{Biên quyết định của chín mô hình trên dữ liệu \texttt{make\_moons} (400 điểm, nhiễu 0,28).}{ch2-boundaries}

% =====================================================================
\section{Học không giám sát}
\label{sec:ml-unsup}

\subsection{Phân cụm K-means}
K-means~\cite{ch2:lloyd1982} chia $n$ điểm thành $k$ cụm bằng cách cực tiểu tổng bình phương khoảng cách tới tâm cụm
(inertia) $J=\sum_{i}\lVert\mathbf{x}_i-\boldsymbol\mu_{c(i)}\rVert^2$. Thuật toán Lloyd luân phiên hai bước cho tới khi
hội tụ: \emph{gán} mỗi điểm vào tâm gần nhất, rồi \emph{cập nhật} mỗi tâm bằng trung bình các điểm của nó
(Hình~\ref{fig:ch2-kmeans}). Mỗi bước không làm tăng $J$ nên thuật toán luôn hội tụ, nhưng chỉ tới cực tiểu cục bộ; vì vậy
thực tế chạy nhiều lần khởi tạo (k-means++, \texttt{n\_init}). Số cụm $k$ được chọn bằng phương pháp ``khuỷu tay'' trên
đường inertia và hệ số silhouette~\cite{ch2:rousseeuw1987} $s=(b-a)/\max(a,b)$.

\hinhpdf{ch2_kmeans_iter}{1}{Thuật toán Lloyd trên 4 cụm: khởi tạo ngẫu nhiên rồi ba vòng gán -- cập nhật tâm (dấu X).}{ch2-kmeans}
\maNguon{ch2_unsup__kmeans_numpy}{K-means (thuật toán Lloyd) cài đặt bằng NumPy}

\subsection{Giảm chiều: PCA}
PCA~\cite{ch2:pearson1901} tìm các hướng trực giao có phương sai lớn nhất. Sau khi tâm hóa dữ liệu, phân rã SVD
$Z=U\Sigma V^\top$ cho các thành phần chính là các hàng của $V^\top$, và tỷ lệ phương sai giải thích của thành phần $k$
là $\sigma_k^2/\sum_j\sigma_j^2$. Chiếu lên $m<d$ thành phần đầu cho biểu diễn ít chiều giữ được nhiều thông tin nhất theo
nghĩa bình phương tối thiểu (Hình~\ref{fig:ch2-pca}).

\hinhpdf{ch2_pca_toy}{1}{PCA trên dữ liệu 2 chiều tương quan: PC1 nằm dọc hướng biến thiên lớn nhất.}{ch2-pca}
\maNguon{ch2_unsup__pca_svd}{PCA bằng phân rã SVD}

\subsection{Thực nghiệm: K-means và PCA trên hồ sơ tiểu đường}
\label{sec:ml-unsup-exp}
Áp dụng lên \SL{ch2_unsup/n} hồ sơ (bốn thuộc tính liên tục đã chuẩn hóa: tuổi, BMI, HbA1c, glucose), đường inertia
giảm đều và silhouette (Hình~\ref{fig:ch2-elbow}) đạt khoảng \SL{ch2_unsup/scan/silhouette/2} ở $k=4$, một giá trị khiêm
tốn cho thấy dữ liệu không có cụm tách bạch rõ. Hai thành phần chính đầu giải thích
\SL{ch2_unsup/pca/explained/0Pct}\% và \SL{ch2_unsup/pca/explained/1Pct}\% phương sai.

\hinhpdf{ch2_unsup_elbow}{1}{Chọn số cụm: inertia (khuỷu tay) và hệ số silhouette theo $k$.}{ch2-elbow}
\hinhpdf{ch2_unsup_pca}{1}{Chiếu hồ sơ lên hai thành phần chính: tô theo cụm K-means (trái) và theo nhãn tiểu đường thật
(phải). Ca bệnh tập trung ở một phía của PC1.}{ch2-unsuppca}
\bang{ch2_unsup_profiles}{Hồ sơ trung bình của bốn cụm K-means ($k=4$) và tỷ lệ tiểu đường trong từng cụm}{ch2-unsup-prof}
\hinhpdf{ch2_unsup_profiles}{1}{Tỷ lệ tiểu đường trong bốn cụm tìm được \emph{không cần nhãn}.}{ch2-unsupprof}

\nhanxet{Dù K-means không hề nhìn thấy nhãn, các cụm có tỷ lệ tiểu đường chênh nhau rõ rệt (Bảng~\ref{tab:ch2-unsup-prof}):
cụm có HbA1c và glucose trung bình cao nhất cũng có tỷ lệ bệnh cao nhất. Học không giám sát vì vậy hữu ích để phân nhóm
nguy cơ và khám phá dữ liệu, nhưng không thay được bộ phân loại có giám sát: silhouette thấp cho thấy ranh giới cụm chỉ là
quy ước, không phải cấu trúc tự nhiên.}

% =====================================================================
\section{Mạng nơ-ron nhân tạo (nhập môn)}
\label{sec:ml-mlp}
Mạng nơ-ron truyền thẳng nhiều lớp (MLP) xếp chồng các phép biến đổi affine và hàm kích hoạt phi tuyến:
$\mathbf{h}^{(l)}=\phi\big(W^{(l)}\mathbf{h}^{(l-1)}+\mathbf{b}^{(l)}\big)$ với $\phi=\mathrm{ReLU}(z)=\max(0,z)$. Hồi quy
logistic chính là MLP không có lớp ẩn. Định lý xấp xỉ phổ quát nói rằng một lớp ẩn đủ rộng xấp xỉ được mọi hàm liên tục
trên miền compact. Tham số được học bằng lan truyền ngược~\cite{ch2:rumelhart1986}: áp dụng quy tắc dây chuyền từ lớp ra
về lớp vào, $\partial\mathcal{L}/\partial W^{(l)}=\boldsymbol\delta^{(l)}\mathbf{h}^{(l-1)\top}$ với
$\boldsymbol\delta^{(l)}=\big(W^{(l+1)\top}\boldsymbol\delta^{(l+1)}\big)\odot\phi'(\mathbf{z}^{(l)})$.

MLP dùng trong TN1 và TN2 (Hình~\ref{fig:ch2-mlp}) có ba lớp ẩn 128--64--32, mỗi lớp gồm Linear, Batch Normalization%
~\cite{ch2:ioffe2015}, ReLU và Dropout, huấn luyện bằng Adam với dừng sớm theo mất mát Validation trên GPU bằng
PyTorch~\cite{ch2:paszke2019}. Chương~3 và~4 sẽ mở rộng ý tưởng này thành mạng tích chập và mạng hồi quy.

\hinhtikz{ch2_mlp}{MLP cho dữ liệu bảng: \SL{ch2_diabetes/split/n_features} đặc trưng vào (TN1), ba lớp ẩn 128--64--32 (Linear, BatchNorm, ReLU, Dropout) và một
nơ-ron ra với sigmoid (phân loại) hoặc tuyến tính (hồi quy).}{ch2-mlp}
\maNguon{ch2_mlp__TabularMLP}{Kiến trúc MLP cho dữ liệu bảng bằng PyTorch}
\maNguon{ch2_mlp__train_epoch}{Một epoch huấn luyện: lan truyền thuận, lan truyền ngược, cập nhật Adam}
\maDai{ch2_mlp__fit_mlp}{Huấn luyện MLP với dừng sớm theo mất mát Validation và giảm tốc độ học}

% =====================================================================
\section{Xử lý và biến đổi đặc trưng}
\label{sec:ml-feature}
\begin{itemize}
  \item \textbf{Chuẩn hóa z-score} $x'=(x-\mu)/\sigma$ (StandardScaler) cho mô hình dựa trên khoảng cách hay gradient
        (KNN, SVM, logistic, MLP); mô hình cây không cần.
  \item \textbf{Mã hóa One-Hot} cho biến phân loại: một hạng mục thành một cột nhị phân. Với mô hình tuyến tính có hệ số
        chặn, bỏ một cột mỗi biến (\texttt{drop="first"}) để tránh đa cộng tuyến hoàn hảo.
  \item \textbf{TF-IDF} cho văn bản~\cite{ch2:salton1988}: trọng số của từ $t$ trong tài liệu $d$ là
        $\text{tf}(t,d)\cdot\log\frac{1+N}{1+\text{df}(t)}+1$, tức từ xuất hiện nhiều trong tài liệu nhưng hiếm trong kho
        thì có trọng số cao.
  \item \textbf{Xử lý mất cân bằng}: trọng số lớp (\texttt{class\_weight="balanced"}, \texttt{pos\_weight}), lấy mẫu
        lại, hoặc đơn giản và hiệu quả nhất là \emph{chọn ngưỡng} trên Validation.
  \item \textbf{Ngoại lai}: chỉ loại khi chắc chắn là lỗi đo/nhập; giá trị cực trị có thật (đường huyết rất cao) mang
        thông tin quan trọng.
\end{itemize}

\maNguon{ch2_data__diabetes_preprocessor}{Bộ tiền xử lý dùng chung: StandardScaler cho thuộc tính số, One-Hot cho thuộc tính phân loại}
\maNguon{ch2_data__split_3way}{Chia Train/Validation/Test 70/15/15 có phân tầng}

% =====================================================================
\section{Cài đặt từ đầu và đối chiếu với thư viện}
\label{sec:ml-scratch}
Để chắc chắn công thức ở các mục trên là đúng chứ không chỉ ``gọi thư viện'', Bảng~\ref{tab:ch2-scratch} đối chiếu các
cài đặt NumPy với scikit-learn trên chính dữ liệu thực nghiệm. Phương trình chuẩn cho hệ số trùng với
\texttt{LinearRegression} tới sai khác lớn nhất \SL{ch2_scratch/linreg/max_coef_diff}; hồi quy logistic tự viết đạt
ROC-AUC \SL{ch2_scratch/logreg/auc_numpy} so với \SL{ch2_scratch/logreg/auc_sklearn} của thư viện, hệ số tương quan
giữa hai vector trọng số \SL{ch2_scratch/logreg/corr_coef}. Khác biệt duy nhất là tốc độ: mini-batch GD bằng Python mất
\SL{ch2_scratch/logreg/time_numpy_s} giây so với \SL{ch2_scratch/logreg/time_sklearn_s} giây của bộ giải L-BFGS.

\bang{ch2_scratch}{Đối chiếu cài đặt NumPy từ đầu với scikit-learn trên tập Test}{ch2-scratch}
\hinhpdf{ch2_scratch_loss}{1}{Mất mát giảm dần của hai thuật toán Gradient Descent tự viết.}{ch2-scratchloss}

% =====================================================================
\section{Thực nghiệm 1: Dự đoán bệnh tiểu đường}
\label{sec:ml-diab}

\subsection{Mô tả bài toán}
Tiểu đường type 2 thường tiến triển âm thầm nhiều năm trước khi được chẩn đoán. Một mô hình sàng lọc dựa trên các chỉ số
sẵn có trong hồ sơ khám (tuổi, BMI, HbA1c, đường huyết, tiền sử tăng huyết áp và tim mạch) giúp khoanh vùng người có nguy
cơ cao để chỉ định xét nghiệm chuyên sâu. Bài toán là \textbf{phân loại nhị phân}: đầu vào là một hồ sơ $\mathbf{x}$, đầu
ra là xác suất $P(\text{diabetes}=1\mid\mathbf{x})$ và nhãn sau khi so với ngưỡng. Vì là bài toán sàng lọc y tế, \emph{bỏ
sót} người bệnh (FN) nghiêm trọng hơn \emph{báo động nhầm} (FP), nên ngoài accuracy ta theo dõi Recall, F1 và PR-AUC.

\subsection{Giới thiệu dữ liệu}
Bộ dữ liệu \emph{Diabetes Prediction Dataset}~\cite{ch2:ds_diabetes} gồm \SL{ch2_diabetes/audit_raw/n_rows} hồ sơ và
\SL{ch2_diabetes/audit_raw/cols} cột: hai biến phân loại (\texttt{gender}, \texttt{smoking\_history}), hai biến nhị phân
(\texttt{hypertension}, \texttt{heart\_disease}), bốn biến liên tục (\texttt{age}, \texttt{bmi}, \texttt{HbA1c\_level},
\texttt{blood\_glucose\_level}) và nhãn \texttt{diabetes}. HbA1c phản ánh đường huyết trung bình 2--3 tháng; theo tiêu
chuẩn ADA~\cite{ch2:ada2023}, HbA1c $\ge 6{,}5\%$ hoặc glucose lúc đói $\ge 126$ mg/dL là ngưỡng chẩn đoán.
Đoạn mã~\ref{code:ch2_data__download} tải ba bộ dữ liệu của chương về \texttt{data/ch2}.

\maNguon{ch2_data__download}{Tải ba bộ dữ liệu của Chương 2 về thư mục dữ liệu cục bộ}

\subsection{Khảo sát dữ liệu}
\bang{ch2_diab_describe}{Thống kê mô tả các thuộc tính số của dữ liệu tiểu đường (trước làm sạch)}{ch2-diab-describe}

Bảng~\ref{tab:ch2-diab-describe} và Hình~\ref{fig:ch2-diab-target} cho thấy ba đặc điểm. Thứ nhất, dữ liệu \textbf{mất
cân bằng}: chỉ \SL{ch2_diabetes/survey/pos_ratePct}\% hồ sơ có bệnh. Thứ hai, không có giá trị thiếu dạng
\texttt{NaN} (\SL{ch2_diabetes/audit_raw/missing} ô), nhưng \texttt{smoking\_history} có hạng mục ``No Info'', về bản chất
là thiếu thông tin được mã hóa thành một nhóm. Thứ ba, có \SL{ch2_diabetes/audit_raw/duplicates} bản ghi trùng lặp hoàn toàn.

\maNguon{ch2_diabetes__survey}{Khảo sát nhanh: kích thước, kiểu dữ liệu, phân bố nhãn và hạng mục}
\maNguon{ch2_data__audit}{Kiểm tra chất lượng: giá trị thiếu, bản ghi trùng và ngoại lai theo quy tắc IQR}
\hinhpdf{ch2_diab_target}{1}{Phân bố nhãn \texttt{diabetes}: lớp dương chiếm chưa tới một phần mười.}{ch2-diab-target}

\subsection{Làm sạch dữ liệu}
Quy tắc IQR đánh dấu một điểm là ngoại lai nếu nằm ngoài $[Q_1-1{,}5\,\mathrm{IQR},\;Q_3+1{,}5\,\mathrm{IQR}]$.
Bảng~\ref{tab:ch2-diab-outliers} cho thấy BMI có \SL{ch2_diabetes/audit_raw/outliers/bmi/n} điểm vượt ngưỡng,
glucose có \SL{ch2_diabetes/audit_raw/outliers/blood_glucose_level/n} điểm và HbA1c có
\SL{ch2_diabetes/audit_raw/outliers/HbA1c_level/n} điểm. Tuy nhiên, Hình~\ref{fig:ch2-diab-outliers} cho thấy các điểm
``ngoại lai'' của glucose và HbA1c tập trung gần như hoàn toàn ở lớp có bệnh: đó chính là tín hiệu cần học, không phải lỗi.

\bang{ch2_diab_outliers}{Ngoại lai theo quy tắc IQR trên bốn thuộc tính liên tục}{ch2-diab-outliers}
\hinhpdf{ch2_diab_outliers}{1}{Biểu đồ hộp theo nhãn của ba thuộc tính có ngoại lai (age không có điểm nào nằm ngoài cận
IQR): giá trị cao của HbA1c và glucose gần như chỉ thuộc người bệnh.}{ch2-diab-outliers}

Vì vậy chiến lược làm sạch (Đoạn mã~\ref{code:ch2_data__clean_diabetes}) là: (1) loại bản ghi trùng; (2) chỉ loại giá trị
\emph{bất khả về sinh lý} (BMI, HbA1c, glucose ngoài khoảng sống được), \emph{giữ nguyên} các giá trị cao có thật;
(3) giữ ``No Info'' như một hạng mục riêng vì việc thiếu thông tin hút thuốc tự nó mang tín hiệu. Sau làm sạch còn
\SL{ch2_diabetes/clean/n_rows} hồ sơ (loại \SL{ch2_diabetes/clean/removed}), trong đó \SL{ch2_diabetes/clean/pos} ca bệnh
(\SL{ch2_diabetes/clean/pos_ratePct}\%).

\maNguon{ch2_data__clean_diabetes}{Làm sạch dữ liệu tiểu đường: bỏ trùng lặp, chỉ loại giá trị bất khả về sinh lý}

\subsection{Biểu diễn dữ liệu}
Sáu thuộc tính số được chuẩn hóa z-score; \texttt{gender} (3 hạng mục) và \texttt{smoking\_history} (6 hạng mục) được mã
hóa One-Hot, tổng cộng \SL{ch2_diabetes/split/n_features} đặc trưng (Đoạn mã~\ref{code:ch2_data__diabetes_preprocessor}).
Bộ tiền xử lý là một \texttt{ColumnTransformer} được fit trên Train và lưu cùng mô hình, nên Web App ở mục
\ref{sec:ml-diab-deploy} biến đổi dữ liệu người dùng nhập vào \emph{đúng như} lúc huấn luyện.

\maNguon{ch2_data__feature_names}{Lấy tên các đặc trưng sau biến đổi để diễn giải hệ số và độ quan trọng}

\subsection{Phân tích khám phá (EDA)}
\hinhpdf{ch2_diab_hist}{1}{Phân bố bốn thuộc tính liên tục theo nhãn (mật độ chuẩn hóa).}{ch2-diab-hist}
\bonhanxet{Phân bố HbA1c và glucose của hai lớp gần như tách rời: người không bệnh tập trung ở HbA1c dưới 6{,}5\% và
glucose dưới 200, còn người bệnh trải rộng sang phải (Hình~\ref{fig:ch2-diab-hist}). Tuổi và BMI của người bệnh dịch
sang phải nhưng chồng lấn lớn.}
{HbA1c và glucose là các xét nghiệm dùng trực tiếp để chẩn đoán, nên tách lớp tốt là đúng với kiến thức y khoa. Dữ liệu
có vẻ được sinh với những ngưỡng sắc nét (glucose nhận ít giá trị rời rạc), điều sẽ có lợi cho mô hình cây.}
{Một mô hình chỉ cần vài ngưỡng trên HbA1c và glucose đã đạt hiệu năng cao; phần khó là các ca bệnh có chỉ số xét nghiệm
``bình thường'', nơi tuổi, BMI và bệnh nền phải bù vào.}

\hinhpdf{ch2_diab_violin}{1}{Biểu đồ violin của bốn thuộc tính liên tục theo nhãn (0: không bệnh, 1: có bệnh).}{ch2-diab-violin}
\hinhpdf{ch2_diab_corr}{1}{Ma trận tương quan Pearson giữa các thuộc tính số và nhãn.}{ch2-diab-corr}

Hình~\ref{fig:ch2-diab-corr} định lượng quan sát trên: tương quan với nhãn cao nhất thuộc về glucose
(\SL{ch2_diabetes/clean/corr_target/blood_glucose_level}) và HbA1c (\SL{ch2_diabetes/clean/corr_target/HbA1c_level}),
tiếp theo là tuổi (\SL{ch2_diabetes/clean/corr_target/age}) và BMI (\SL{ch2_diabetes/clean/corr_target/bmi}). Các thuộc
tính đầu vào tương quan với nhau ở mức thấp, nên mô hình tuyến tính không gặp vấn đề đa cộng tuyến nghiêm trọng.

\hinhpdf{ch2_diab_rates}{1}{Tỷ lệ tiểu đường theo giới tính, tiền sử hút thuốc, tăng huyết áp và bệnh tim.}{ch2-diab-rates}
\bonhanxet{Tỷ lệ bệnh ở người tăng huyết áp là \SL{ch2_diabetes/clean/rate_hyper/1Pct}\% so với
\SL{ch2_diabetes/clean/rate_hyper/0Pct}\% ở người không tăng huyết áp; với bệnh tim là
\SL{ch2_diabetes/clean/rate_heart/1Pct}\% so với \SL{ch2_diabetes/clean/rate_heart/0Pct}\%. Nhóm từng hút thuốc
(\texttt{former}) có tỷ lệ \SL{ch2_diabetes/clean/rate_by_smoking/formerPct}\%, cao hơn nhóm chưa từng hút
(\SL{ch2_diabetes/clean/rate_by_smoking/neverPct}\%).}
{Tăng huyết áp, bệnh tim và tiểu đường là các bệnh đồng mắc của hội chứng chuyển hóa. Nhóm \texttt{former} thường lớn tuổi
hơn, nên một phần chênh lệch đến từ tuổi chứ không phải từ hút thuốc.}
{Các biến nhị phân có ích nhưng yếu hơn xét nghiệm; tương quan không phải nhân quả, nên không diễn giải hệ số của
\texttt{smoking\_history} như ``tác động'' của hút thuốc.}

\hinhpdf{ch2_diab_scatter}{1}{Biểu đồ phân tán HbA1c theo glucose (mẫu ngẫu nhiên): ca bệnh nằm ở phần trên bên phải và
dọc theo các ngưỡng.}{ch2-diab-scatter}

\subsection{Chia tập dữ liệu}
Dữ liệu sạch được chia phân tầng 70/15/15 (Đoạn mã~\ref{code:ch2_data__split_3way}). Bảng~\ref{tab:ch2-diab-split} xác
nhận tỷ lệ lớp dương của ba tập gần như bằng nhau, nên điểm số trên Validation và Test có thể so sánh trực tiếp.
\bang{ch2_diab_split}{Kích thước và tỷ lệ lớp dương của ba tập}{ch2-diab-split}

\subsection{Xây dựng mô hình}
Chín họ mô hình được huấn luyện trên cùng tập Train, cùng bộ tiền xử lý (Đoạn mã~\ref{code:ch2_models__classifiers}):
\begin{itemize}
  \item \textbf{Baseline}: luôn dự đoán lớp đa số; mức sàn để biết một mô hình có ``học'' được gì không.
  \item \textbf{Logistic Regression} với \texttt{class\_weight="balanced"} để bù mất cân bằng.
  \item \textbf{KNN} $k=25$ trọng số theo khoảng cách; \textbf{Gaussian Naive Bayes}.
  \item \textbf{Decision Tree} độ sâu 8, tối thiểu 20 mẫu ở lá; \textbf{Random Forest} 300 cây.
  \item \textbf{XGBoost} và \textbf{LightGBM} 400 cây, \texttt{scale\_pos\_weight} bằng tỷ lệ âm/dương
        (\SL{ch2_diabetes/pos_weight}).
  \item \textbf{SVM RBF}: huấn luyện trên \SL{ch2_diabetes/models/svm/n_train} mẫu con do độ phức tạp $O(n^2)$ của bộ giải.
  \item \textbf{MLP} PyTorch (mục~\ref{sec:ml-mlp}) với \SL{ch2_diabetes/models/mlp/n_params} tham số, chạy trên GPU.
\end{itemize}

\maNguon{ch2_models__classifiers}{Khai báo chín bộ phân loại với siêu tham số dùng trong thực nghiệm}
\maNguon{ch2_models__clf_pipeline}{Ghép bộ tiền xử lý và mô hình thành một Pipeline duy nhất}
\maNguon{ch2_diabetes__fit_classic}{Huấn luyện một mô hình cổ điển, đo thời gian và lưu ra đĩa}

\subsection{Đánh giá mô hình}
\bang{ch2_diab_val}{Kết quả trên tập Validation ở ngưỡng tối ưu F1 của từng mô hình}{ch2-diab-val}
\hinhpdf{ch2_diab_compare}{1}{So sánh Precision, Recall, F1 và ROC-AUC của các mô hình trên Validation (ngưỡng tối ưu F1).}{ch2-diab-compare}

\bonhanxet{Baseline đạt accuracy \SL{ch2_diabetes/val/baseline/accPct}\% nhưng F1 bằng \SL{ch2_diabetes/val/baseline/f1}:
accuracy hoàn toàn không phản ánh khả năng phát hiện bệnh. Các mô hình cây tổ hợp và MLP đứng đầu: XGBoost đạt F1
\SL{ch2_diabetes/val/xgb/f1} và ROC-AUC \SL{ch2_diabetes/val/xgb/auc}, MLP đạt F1 \SL{ch2_diabetes/val/mlp/f1}, LightGBM
\SL{ch2_diabetes/val/lgbm/f1}. Chỉ một cây quyết định đơn đã đạt F1 \SL{ch2_diabetes/val/tree/f1} với Precision rất cao
\SL{ch2_diabetes/val/tree/prec}. Logistic Regression (\SL{ch2_diabetes/val/logreg/f1}) và KNN
(\SL{ch2_diabetes/val/knn/f1}) kém hơn rõ; Naive Bayes thấp nhất (\SL{ch2_diabetes/val/nb/f1}).}
{Quan hệ giữa xét nghiệm và nhãn có dạng ngưỡng (đã thấy ở EDA), đúng sở trường của cây: một phép tách
``HbA1c $>$ \SL{ch2_tree/nodes/n0/thr}'' đã cho lá thuần. Mô hình tuyến tính phải xấp xỉ một bậc thang bằng một mặt phẳng nên mất độ chính xác ở
vùng giáp ranh. Naive Bayes giả định các thuộc tính độc lập và có phân phối chuẩn, trong khi HbA1c và glucose rời rạc,
lệch phải và cùng phụ thuộc vào bệnh.}
{Với dữ liệu bảng có quan hệ ngưỡng, gradient boosting là lựa chọn mặc định đáng tin cậy. MLP gần bằng nhưng tốn hơn
nhiều thời gian huấn luyện (\SL{ch2_diabetes/models/mlp/train_s} giây so với \SL{ch2_diabetes/models/xgb/train_s} giây).}

\subsubsection{Ảnh hưởng của ngưỡng quyết định}
\bang{ch2_diab_thr}{Precision, Recall, F1 ở ngưỡng 0,5 và ở ngưỡng tối ưu F1 (*) chọn trên Validation}{ch2-diab-thr}
\hinhpdf{ch2_diab_threshold}{1}{Trái: F1 trên Validation theo ngưỡng của năm mô hình (chấm tròn là ngưỡng tối ưu). Phải: Precision và
Recall của XGBoost theo ngưỡng; đường dọc là ngưỡng đã chọn.}{ch2-diab-thr}

Bảng~\ref{tab:ch2-diab-thr} cho thấy chọn ngưỡng giúp nhiều nhất cho các mô hình có xác suất \emph{lệch chuẩn}: Naive Bayes
và SVM. Với các mô hình đã cân bằng lớp như XGBoost, ngưỡng tối ưu (\SL{ch2_diabetes/thresholds/xgb}) gần 0,5. Hình
\ref{fig:ch2-diab-thr} là lựa chọn thực tế: hạ ngưỡng xuống tăng Recall, đổi lại Precision giảm nhanh. Nếu mục tiêu sàng lọc
đòi Recall cao hơn, ta chọn ngưỡng thấp hơn ngưỡng tối ưu F1 và chấp nhận nhiều ca cần xét nghiệm lại hơn.

\subsubsection{Đường ROC, Precision--Recall và ma trận nhầm lẫn}
\hinhpdf{ch2_diab_roc_pr}{1}{Đường ROC (trái) và Precision--Recall (phải) của các mô hình trên tập Test; chú giải ghi AUC và AP.}{ch2-diab-rocpr}
\hinhpdf{ch2_diab_confusion}{1}{Ma trận nhầm lẫn của chín mô hình trên tập Test ở ngưỡng chọn trên Validation (tỷ lệ theo hàng, số
lượng trong ngoặc).}{ch2-diab-cm}

Trên đường ROC, mọi mô hình trừ Naive Bayes đều ôm sát góc trên bên trái nên khó phân biệt. Đường Precision--Recall
phân biệt rõ hơn: mức ngẫu nhiên là \SL{ch2_diabetes/test/baseline/ap} (tỷ lệ lớp dương), XGBoost đạt PR-AUC
\SL{ch2_diabetes/test/xgb/ap} còn Logistic Regression \SL{ch2_diabetes/test/logreg/ap}. Ma trận nhầm lẫn
(Hình~\ref{fig:ch2-diab-cm}) cho thấy điểm yếu chung của các mô hình tốt nhất: chúng gần như không báo động nhầm nhưng bỏ sót khoảng
một phần ba số ca bệnh, là những ca có HbA1c và glucose nằm trong vùng bình thường.

\subsubsection{Độ phức tạp, đường học và kiểm định chéo}
\haihinh{ch2_diab_depth}{Decision Tree theo độ sâu}{ch2_diab_knnk}{KNN theo số láng giềng $k$}{Đường kiểm định: ROC-AUC
trên Train và Validation theo siêu tham số độ phức tạp.}{ch2-diab-curves}

Hình~\ref{fig:ch2-diab-curves}a minh họa đúng lý thuyết thiên lệch -- phương sai: AUC Train tăng đơn điệu theo độ sâu,
còn AUC Validation đạt đỉnh \SL{ch2_diabetes/curve_depth/val_auc/7} ở độ sâu \SL{ch2_diabetes/curve_depth/depth/7} rồi giảm
về \SL{ch2_diabetes/curve_depth/val_auc/11} ở độ sâu \SL{ch2_diabetes/curve_depth/depth/11} khi cây học thuộc nhiễu. KNN
cho chiều ngược lại: $k$ nhỏ quá khớp, $k$ lớn làm mượt và AUC Validation tăng dần.

\hinhpdf{ch2_diab_learning}{1}{Đường học của XGBoost: F1 trên Train và Validation theo số mẫu huấn luyện.}{ch2-diab-learning}
\bang{ch2_diab_cv}{Kiểm định chéo 5-fold trên tập Train: trung bình $\pm$ độ lệch chuẩn}{ch2-diab-cv}
\hinhpdf{ch2_diab_cv}{1}{Phân bố điểm số 5-fold của bốn mô hình.}{ch2-diab-cvfig}

Đường học (Hình~\ref{fig:ch2-diab-learning}) cho thấy khoảng cách Train -- Validation thu hẹp khi tăng dữ liệu và
Validation gần như bão hòa từ khoảng nửa tập Train: thêm dữ liệu cùng loại sẽ không cải thiện đáng kể, muốn tốt hơn cần
thêm \emph{thông tin mới} (đặc trưng khác). Kiểm định chéo (Bảng~\ref{tab:ch2-diab-cv}) xác nhận thứ hạng ổn định:
XGBoost có ROC-AUC $\SL{ch2_diabetes/cv/xgb/test_roc_auc/mean}\pm\SL{ch2_diabetes/cv/xgb/test_roc_auc/std}$, cao nhất và
dao động nhỏ nhất giữa các fold.

\subsubsection{Diễn giải mô hình}
\hinhpdf{ch2_diab_importance}{1}{Độ quan trọng đặc trưng: hoán vị (permutation, mức giảm ROC-AUC trên Validation) và gain của XGBoost (10
đặc trưng đầu).}{ch2-diab-imp}
\hinhpdf{ch2_diab_coef}{1}{Hệ số của Logistic Regression trên đặc trưng đã chuẩn hóa.}{ch2-diab-coef}

\bonhanxet{Cả ba cách đo độ quan trọng (hoán vị, gain của XGBoost và mức giảm Gini của Random Forest, với HbA1c
\SL{ch2_diabetes/importance/rf/HbA1c_level}) đều xếp HbA1c và glucose ở hai vị trí đầu, cách xa phần còn lại: hoán vị HbA1c làm
ROC-AUC của XGBoost giảm \SL{ch2_diabetes/importance/perm/HbA1c_level/mean}, glucose
\SL{ch2_diabetes/importance/perm/blood_glucose_level/mean}, tuổi chỉ \SL{ch2_diabetes/importance/perm/age/mean}. Hệ số
Logistic Regression lớn nhất cũng thuộc HbA1c (\SL{ch2_diabetes/importance/logreg_coef/HbA1c_level}).}
{Mô hình học đúng các yếu tố mà y khoa dùng để chẩn đoán, không phải một quy luật giả tạo từ dữ liệu. Hệ số dương lớn của
tuổi và BMI phù hợp với hiểu biết về yếu tố nguy cơ.}
{Mô hình đủ đáng tin để làm công cụ sàng lọc; nhưng chính vì phụ thuộc mạnh vào xét nghiệm, nó không giúp được nhiều cho
người chưa từng đo HbA1c. Một mô hình ``không xét nghiệm'' chỉ dùng tuổi, BMI, bệnh nền sẽ là bài toán khác, khó hơn.}

\hinhpdf{ch2_diab_mlp_calib}{1}{Trái: mất mát Train/Validation của MLP theo epoch (dừng sớm). Phải: biểu đồ hiệu chuẩn
xác suất của các mô hình trên tập Test.}{ch2-diab-mlp}
Biểu đồ hiệu chuẩn (Hình~\ref{fig:ch2-diab-mlp}) kiểm tra xác suất dự đoán có ``trung thực'' hay không~\cite{ch2:niculescu2005}:
trong các hồ sơ được gán xác suất 0,7, có thực sự khoảng 70\% bị bệnh? Logistic Regression, XGBoost và MLP bám sát
đường chéo, tức xác suất của chúng dùng được như mức nguy cơ. Naive Bayes thì ngược lại: nó gán xác suất gần 1 cho rất
nhiều hồ sơ mà tỷ lệ bệnh thật trong nhóm đó thấp hơn nhiều, hệ quả của giả định độc lập bị vi phạm. Đây là lý do ngưỡng
tối ưu của Naive Bayes lên tới \SL{ch2_diabetes/thresholds/nb}, và lý do ngưỡng phải được chọn trên Validation cho từng
mô hình thay vì mặc định 0,5. Web App vì vậy hiển thị xác suất như một ``mức nguy cơ'' để tham khảo, không phải chẩn đoán.

\subsection{Lựa chọn mô hình và kết quả trên tập Test}
Mô hình được chọn theo F1 trên Validation, tiêu chí phụ là ROC-AUC và độ ổn định qua kiểm định chéo: \textbf{XGBoost}.
Tập Test chỉ được dùng một lần sau khi đã cố định mô hình và ngưỡng.

\bang{ch2_diab_test}{Kết quả trên tập Test (ngưỡng cố định từ Validation)}{ch2-diab-test}
\bang{ch2_diab_best_cm}{Ma trận nhầm lẫn của XGBoost trên tập Test}{ch2-diab-bestcm}

\begin{ketqua}{Kết quả cuối cùng -- Thực nghiệm 1}
Mô hình: \SL{ch2_diabetes/best/label} (ngưỡng \SL{ch2_diabetes/best/thr})\\
Accuracy = \SL{ch2_diabetes/best/acc} · Precision = \SL{ch2_diabetes/best/prec} · Recall = \SL{ch2_diabetes/best/rec}\\
F1 = \SL{ch2_diabetes/best/f1} · ROC-AUC = \SL{ch2_diabetes/best/auc} · PR-AUC = \SL{ch2_diabetes/best/ap}\\
Bỏ sót \SL{ch2_diabetes/test/xgb/fn} ca, báo động nhầm \SL{ch2_diabetes/test/xgb/fp} ca trên \SL{ch2_diabetes/split/test} hồ sơ Test
\end{ketqua}

Kết quả Test (F1 \SL{ch2_diabetes/test/xgb/f1}) rất sát Validation (\SL{ch2_diabetes/val/xgb/f1}), xác nhận không có quá
khớp vào Validation khi chọn mô hình và ngưỡng. Precision \SL{ch2_diabetes/best/prec} nghĩa là gần như mọi cảnh báo đều
đúng; Recall \SL{ch2_diabetes/best/rec} là giới hạn chính: khoảng ba phần mười ca bệnh có chỉ số xét nghiệm không khác người
khỏe, và không mô hình nào trong chín mô hình phát hiện được chúng từ dữ liệu này.

\subsection{Triển khai hệ thống}
\label{sec:ml-diab-deploy}
Mô hình được triển khai thành dịch vụ dự đoán trong Web App chung của tiểu luận (Hình~\ref{fig:ch2-deploy}). Mọi mô hình
đã huấn luyện được lưu bằng \texttt{joblib} (MLP lưu \texttt{state\_dict} của PyTorch) cùng bộ tiền xử lý. Module
\texttt{tl3/serve\_ch2.py} cung cấp một hàm duy nhất \texttt{predict(exp\_id, payload, model\_id)}; Web App gọi hàm này
qua endpoint \texttt{POST /api/predict} với thân JSON
\texttt{\{"exp\_id": "ch2\_diabetes", "model": "xgb", "payload": \{...\}\}}. Mô hình được nạp lười một lần rồi giữ trong
bộ nhớ (\texttt{lru\_cache}), nên mỗi lượt dự đoán chỉ mất vài mili giây.

\hinhtikz{ch2_deploy}{Kiến trúc triển khai: trình duyệt và điện thoại gửi JSON tới Web App, Web App gọi
\texttt{serve\_ch2.predict}, mô hình được nạp lười từ \texttt{models/ch2}.}{ch2-deploy}
\maNguon{serve_ch2__predict}{Điểm vào dùng chung của dịch vụ dự đoán Chương 2}
\maNguon{serve_ch2__predict_diabetes}{Dự đoán nguy cơ tiểu đường: tiền xử lý, xác suất, so với ngưỡng đã chọn}

Kết quả trả về gồm \texttt{label} (nhãn tiếng Việt), \texttt{proba} (xác suất của hai lớp), \texttt{details} (ngưỡng và
xác suất lớp dương), tên mô hình và thời gian xử lý. Người dùng có thể đổi mô hình (XGBoost, LightGBM, Random Forest,
Logistic Regression, Decision Tree, KNN, Naive Bayes, MLP) để so sánh trực tiếp cùng một hồ sơ.

\anhUI{ch2_diabetes_web}{Giao diện Web App -- dự đoán nguy cơ tiểu đường trên máy tính.}{ch2-diab-web}
\anhMobile{ch2_diabetes_mobile}{Giao diện dự đoán tiểu đường trên điện thoại (giao diện đáp ứng).}{ch2-diab-mobile}

% =====================================================================
\section{Thực nghiệm 2: Dự đoán giá nhà}
\label{sec:ml-house}

\subsection{Mô tả bài toán}
Định giá bất động sản là bài toán \textbf{hồi quy} kinh điển: từ các thuộc tính của căn nhà (diện tích, số phòng, số tầng,
tuổi nhà, tiện ích, thành phố, mức độ nội thất) dự đoán giá bán. Mô hình định giá tự động hỗ trợ người mua kiểm tra mức giá
rao bán, ngân hàng thẩm định tài sản thế chấp và sàn giao dịch gợi ý giá. Sai số được đo bằng MAE (đơn vị tiền),
RMSE, $R^2$ và MAPE.

\subsection{Giới thiệu và khảo sát dữ liệu}
Bộ dữ liệu \emph{House Price Prediction}~\cite{ch2:ds_house} gồm \SL{ch2_house/audit/n_rows} căn nhà ở bảy thành phố Ấn Độ,
\SL{ch2_house/audit/cols} cột: bảy thuộc tính số (diện tích, phòng ngủ, phòng tắm, số tầng, chỗ đỗ xe, tuổi nhà, điểm khu
vực), tám thuộc tính phân loại và nhãn \texttt{Price}. Dữ liệu sạch: \SL{ch2_house/audit/missing} giá trị thiếu,
\SL{ch2_house/audit/duplicates} bản ghi trùng; chỉ \SL{ch2_house/audit/outliers/Price/n} giá trị giá nằm ngoài khoảng IQR.

\bang{ch2_house_describe}{Thống kê mô tả các thuộc tính số của dữ liệu giá nhà}{ch2-house-describe}
\hinhpdf{ch2_house_price}{1}{Phân bố giá nhà (trái) và giá trung bình theo số phòng ngủ (phải).}{ch2-house-price}

Phân bố giá gần đối xứng (độ lệch \SL{ch2_house/survey/skew_price}), nên không cần biến đổi logarit như thường gặp với dữ
liệu giá thật. Giá trung bình tăng đều theo số phòng ngủ: từ \SL{ch2_house/survey/price_by_bedrooms/1} với một phòng lên
\SL{ch2_house/survey/price_by_bedrooms/5} với năm phòng.

\subsection{Phân tích khám phá}
\hinhpdf{ch2_house_numeric}{1}{Quan hệ giữa giá và các thuộc tính số; đường đỏ là xu hướng tuyến tính.}{ch2-house-num}
\hinhpdf{ch2_house_cats}{1}{Phân bố giá theo thành phố, mức độ nội thất, đối tượng thuê và nguồn nước.}{ch2-house-cats}
\hinhpdf{ch2_house_corr}{1}{Tương quan Pearson giữa các thuộc tính số và giá.}{ch2-house-corr}

\bonhanxet{Số phòng ngủ có tương quan với giá cao nhất (\SL{ch2_house/survey/corr_target/Bedrooms}), sau đó là diện
tích (\SL{ch2_house/survey/corr_target/Area}) và số phòng tắm (\SL{ch2_house/survey/corr_target/Bathrooms}); tuổi nhà
tương quan âm yếu (\SL{ch2_house/survey/corr_target/Age}). Giá trung bình giữa các thành phố chênh không nhiều
(Mumbai \SL{ch2_house/survey/price_by_city/Mumbai}, Chennai \SL{ch2_house/survey/price_by_city/Chennai}), trong khi nhà có
nội thất đầy đủ đắt hơn rõ rệt (\SL{ch2_house/survey/price_by_furnishing/Furnished} so với
\SL{ch2_house/survey/price_by_furnishing/Unfurnished}).}
{Các đám mây điểm trong Hình~\ref{fig:ch2-house-num} có dạng dải thẳng với nhiễu đồng đều, không có đường cong hay
tương tác rõ rệt: dấu hiệu dữ liệu được sinh bởi một quan hệ gần tuyến tính cộng nhiễu. Việc thành phố hầu như không ảnh
hưởng tới giá cũng khác với thị trường thật, nơi Mumbai đắt hơn nhiều lần các thành phố khác.}
{Cần thận trọng khi diễn giải: đây là bộ dữ liệu tổng hợp. Mô hình tuyến tính có thể đã là mô hình ``đúng'', và các mô
hình linh hoạt hơn khó có lợi thế. Thực nghiệm này vì vậy là dịp kiểm chứng điều đó chứ không giả định cây luôn thắng.}

\subsection{Biểu diễn và chia tập}
Thuộc tính số được chuẩn hóa z-score; thuộc tính phân loại được mã hóa One-Hot với \texttt{drop="first"} để hệ số của
mô hình tuyến tính có nghĩa so với một nhóm tham chiếu (ví dụ Bangalore, nhà có nội thất), tổng cộng
\SL{ch2_house/split/n_features} đặc trưng. Dữ liệu chia ngẫu nhiên 70/15/15: \SL{ch2_house/split/train} căn để huấn luyện,
\SL{ch2_house/split/val} kiểm định và \SL{ch2_house/split/test} kiểm thử.

\maNguon{ch2_data__house_preprocessor}{Bộ tiền xử lý cho dữ liệu giá nhà (One-Hot bỏ hạng mục tham chiếu)}

\subsection{Xây dựng mô hình}
Mười mô hình hồi quy được so sánh (Đoạn mã~\ref{code:ch2_models__regressors}): Baseline (dự đoán trung bình), Linear
Regression (OLS), Ridge, Lasso, hồi quy đa thức bậc 2 có Ridge, KNN, Decision Tree, Random Forest, XGBoost, SVR RBF (biến
mục tiêu được chuẩn hóa qua \texttt{TransformedTargetRegressor} vì SVR nhạy với thang của $y$) và MLP. Đoạn mã
\ref{code:ch2_house__regularization_path} tính \emph{đường chính quy hóa}: hệ số của Ridge và Lasso khi $\lambda$ thay đổi.

\maNguon{ch2_models__regressors}{Khai báo mười mô hình hồi quy}
\maNguon{ch2_models__reg_pipeline}{Pipeline tiền xử lý + mô hình hồi quy}
\maNguon{ch2_house__regularization_path}{Đường chính quy hóa của Ridge và Lasso theo $\lambda$}

\subsection{Đánh giá mô hình}
\bang{ch2_house_val}{Kết quả hồi quy trên tập Validation}{ch2-house-val}
\hinhpdf{ch2_house_compare}{1}{$R^2$ và RMSE của các mô hình trên Validation.}{ch2-house-compare}

\bonhanxet{Ba mô hình tuyến tính đứng đầu và gần như bằng nhau: Lasso $R^2=\SL{ch2_house/val/lasso/r2}$, Linear
\SL{ch2_house/val/linear/r2}, Ridge \SL{ch2_house/val/ridge/r2}. MLP (\SL{ch2_house/val/mlp/r2}) và XGBoost
(\SL{ch2_house/val/xgb/r2}) theo sau, còn Random Forest (\SL{ch2_house/val/rf/r2}), KNN (\SL{ch2_house/val/knn/r2}) và
Decision Tree (\SL{ch2_house/val/tree/r2}) kém hơn rõ. Đặc trưng đa thức bậc 2 cũng làm $R^2$ giảm
(\SL{ch2_house/val/poly/r2}).}
{Khi quan hệ thật là tuyến tính, mô hình tuyến tính có thiên lệch bằng 0 và phương sai nhỏ nhất. Cây và KNN phải xấp xỉ một
mặt phẳng bằng các bậc thang hoặc trung bình cục bộ; với chỉ \SL{ch2_house/split/train} mẫu huấn luyện trong không gian
\SL{ch2_house/split/n_features} chiều, phương sai của chúng lớn. Đặc trưng bậc 2 thêm hàng trăm tham số không cần thiết,
đúng kịch bản quá khớp của mục~\ref{sec:ml-overfit}.}
{``Mô hình phức tạp hơn thì tốt hơn'' không phải là quy luật. Trên dữ liệu nhỏ và gần tuyến tính, hồi quy tuyến tính có chính
quy hóa vừa chính xác nhất, vừa nhanh nhất, vừa diễn giải được. Phần $R^2$ còn lại (khoảng một phần tư phương sai) là nhiễu
không mô hình nào giải thích được từ các thuộc tính đã có.}

\hinhpdf{ch2_house_pred}{1}{Giá dự đoán so với giá thật trên tập Test của Lasso và Random Forest; đường chéo là dự đoán hoàn hảo.}{ch2-house-pred}
\hinhpdf{ch2_house_resid}{1}{Phần dư của Lasso: theo giá dự đoán (trái) và phân bố (phải).}{ch2-house-resid}

Phần dư của Lasso (Hình~\ref{fig:ch2-house-resid}) phân bố đối xứng quanh 0 và không có hình phễu theo giá dự đoán, tức
giả định phương sai sai số không đổi của OLS được thỏa mãn. Hình~\ref{fig:ch2-house-pred} cho thấy đám điểm của Lasso bám
đường chéo chặt hơn đám điểm của Random Forest, đúng với chênh lệch $R^2$ giữa hai mô hình trên tập Test.

\subsubsection{Hệ số, đường chính quy hóa và kiểm định chéo}
\hinhpdf{ch2_house_coef}{1}{Hệ số của Linear Regression (đơn vị tiền trên một độ lệch chuẩn của thuộc tính số, hoặc so với
nhóm tham chiếu của thuộc tính phân loại).}{ch2-house-coef}
\hinhpdf{ch2_house_regpath}{1}{Đường chính quy hóa: hệ số Ridge co dần đều về 0, còn hệ số Lasso lần lượt bằng đúng 0 khi
$\lambda$ tăng.}{ch2-house-regpath}

Hệ số tuyến tính cho câu trả lời có thể dùng ngay: tăng số phòng ngủ thêm một độ lệch chuẩn làm giá tăng
\SL{ch2_house/linear_coef/Bedrooms}, nhà không nội thất rẻ hơn nhà có nội thất \SL{ch2_house/linear_coef/Furnishing_Unfurnished}
khi các yếu tố khác như nhau, còn hệ số của các thành phố đều nhỏ. Đường chính quy hóa (Hình~\ref{fig:ch2-house-regpath})
minh họa trực quan khác biệt L1/L2 ở mục~\ref{sec:ml-linreg}: Lasso loại các đặc trưng thành phố đầu tiên, đúng những
đặc trưng ít thông tin nhất.

\bang{ch2_house_cv}{Kiểm định chéo 5-fold ($R^2$ trung bình $\pm$ độ lệch chuẩn, RMSE trung bình)}{ch2-house-cv}
\hinhpdf{ch2_house_cv}{1}{$R^2$ trung bình $\pm$ 2 độ lệch chuẩn qua 5 fold của các mô hình.}{ch2-house-cvfig}
\hinhpdf{ch2_house_learning}{1}{Trái: đường học của Lasso ($R^2$ theo số mẫu huấn luyện). Phải: mất mát Train/Validation của MLP theo epoch.}{ch2-house-learning}

Kiểm định chéo xác nhận thứ hạng: Lasso $R^2=\SL{ch2_house/cv/lasso/r2_mean}\pm\SL{ch2_house/cv/lasso/r2_std}$, XGBoost
$\SL{ch2_house/cv/xgb/r2_mean}\pm\SL{ch2_house/cv/xgb/r2_std}$, Decision Tree dao động mạnh nhất
($\pm\SL{ch2_house/cv/tree/r2_std}$). Đường học của Lasso (Hình~\ref{fig:ch2-house-learning}) hội tụ ngay từ vài trăm mẫu và khoảng cách Train --
Validation rất nhỏ: mô hình đã đạt trần của dữ liệu, thêm mẫu không giúp được. Ngược lại, các mô hình linh hoạt như
Random Forest cần nhiều dữ liệu hơn nhiều để bù phương sai, điều mà \SL{ch2_house/split/train} mẫu không đáp ứng được;
MLP dừng sớm sau \SL{ch2_house/models/mlp/epochs} epoch khi mất mát Validation ngừng giảm.

\subsection{Lựa chọn mô hình và kết quả trên tập Test}
\bang{ch2_house_test}{Kết quả hồi quy trên tập Test}{ch2-house-test}
\bang{ch2_house_quartile}{MAE của Lasso theo bốn khoảng giá (tứ phân vị của giá thật) trên tập Test}{ch2-house-quart}

\begin{ketqua}{Kết quả cuối cùng -- Thực nghiệm 2}
Mô hình: \SL{ch2_house/best/label}\\
$R^2$ = \SL{ch2_house/best/r2} · MAE = \SL{ch2_house/best/mae} · RMSE = \SL{ch2_house/best/rmse} · MAPE = \SL{ch2_house/best/mapePct}\%\\
Baseline (trung bình): $R^2$ = \SL{ch2_house/test/baseline/r2} · MAE = \SL{ch2_house/test/baseline/mae}
\end{ketqua}

Lasso được chọn vì đứng đầu trên cả Validation và kiểm định chéo. Trên Test, nó giảm MAE so với baseline còn chưa tới một
nửa; sai số tương đối trung bình \SL{ch2_house/best/mapePct}\% là mức chấp nhận được cho một công cụ tham khảo giá. Bảng
\ref{tab:ch2-house-quart} cho thấy sai số lớn nhất ở nhóm nhà đắt nhất, nơi có ít mẫu và giá biến thiên nhiều. $R^2$ trên
Test cao hơn Validation chủ yếu do phép chia ngẫu nhiên trên một tập nhỏ; con số đáng tin hơn là trung bình 5-fold.

\subsection{Triển khai hệ thống}
Dịch vụ \texttt{ch2\_house} nhận hồ sơ căn nhà gồm bảy trường số và tám trường lựa chọn, trả về giá dự đoán ở trường
\texttt{value}, kèm đơn vị trong \texttt{details}. Người dùng có thể chuyển giữa Lasso, Linear, Ridge, XGBoost, đa thức,
Random Forest, SVR, KNN và Decision Tree để thấy cùng một căn nhà được định giá khác nhau thế nào.

\maNguon{serve_ch2__predict_house}{Dự đoán giá nhà từ hồ sơ người dùng nhập vào}

\anhUI{ch2_house_web}{Giao diện Web App -- dự đoán giá nhà.}{ch2-house-web}
\anhMobile{ch2_house_mobile}{Giao diện dự đoán giá nhà trên điện thoại.}{ch2-house-mobile}

% =====================================================================
\section{Thực nghiệm 3: Phân loại nhận xét khách hàng}
\label{sec:ml-cmt}

\subsection{Mô tả bài toán}
Doanh nghiệp thương mại điện tử nhận hàng chục nghìn nhận xét mỗi tháng; đọc thủ công là bất khả thi. Bài toán là từ văn bản
nhận xét, dự đoán khách hàng có \textbf{giới thiệu} sản phẩm hay không (phân loại nhị phân, nhãn \texttt{Recommended IND}),
và mở rộng: dự đoán số sao 1--5 (phân loại nhiều lớp). Lớp quan trọng về kinh doanh là lớp \emph{không giới thiệu}, vì đó
là những khách hàng cần được chăm sóc lại. Đây cũng là lớp thiểu số, nên chỉ số chính là \textbf{macro-F1} và F1 của lớp âm.

\subsection{Giới thiệu và khảo sát dữ liệu}
Bộ dữ liệu \emph{Women's E-Commerce Clothing Reviews}~\cite{ch2:ds_reviews} gồm \SL{ch2_comments/survey/raw_rows} nhận
xét, mỗi nhận xét có tiêu đề, nội dung, số sao, nhãn giới thiệu và nhóm sản phẩm. Có \SL{ch2_comments/survey/missing_text}
nhận xét không có nội dung (bị loại) và \SL{ch2_comments/survey/missing_title} nhận xét không có tiêu đề (giữ lại, dùng
nội dung). Sau làm sạch còn \SL{ch2_comments/survey/n_rows} nhận xét, trong đó
\SL{ch2_comments/survey/pos_ratePct}\% được giới thiệu. Văn bản đầu vào là tiêu đề nối với nội dung
(Đoạn mã~\ref{code:ch2_data__load_reviews}).

\maNguon{ch2_data__load_reviews}{Nạp và làm sạch dữ liệu nhận xét: bỏ nhận xét rỗng, ghép tiêu đề với nội dung}
\maNguon{ch2_data__reviews_splits}{Chia phân tầng dữ liệu nhận xét theo nhãn giới thiệu}

\hinhpdf{ch2_cmt_labels}{1}{Trái: phân bố số sao. Phải: tỷ lệ giới thiệu theo số sao.}{ch2-cmt-labels}
\hinhpdf{ch2_cmt_length}{1}{Phân bố độ dài nhận xét (số từ) theo nhãn.}{ch2-cmt-length}
\hinhpdf{ch2_cmt_dept}{1}{Số nhận xét theo nhóm sản phẩm.}{ch2-cmt-dept}

\bonhanxet{Số sao lệch hẳn về 5 sao (\SL{ch2_comments/survey/rating_counts/5} nhận xét). Nhãn giới thiệu gần như trùng với
số sao: tỷ lệ giới thiệu là \SL{ch2_comments/survey/rec_by_rating/1Pct}\% ở 1 sao, \SL{ch2_comments/survey/rec_by_rating/3Pct}\%
ở 3 sao và \SL{ch2_comments/survey/rec_by_rating/5Pct}\% ở 5 sao. Nhận xét có độ dài trung vị
\SL{ch2_comments/survey/words_median} từ, phân vị 95 là \SL{ch2_comments/survey/words_p95} từ, độ dài của hai lớp tương
đương.}
{Nhóm 3 sao là vùng mập mờ: khách vừa khen vừa chê và việc có giới thiệu hay không phụ thuộc vào cảm nhận tổng thể. Độ dài
không phân biệt được hai lớp, nên mô hình phải dựa vào nội dung từ ngữ.}
{Không được dùng số sao làm đặc trưng (gần như là nhãn, gây rò rỉ). Lỗi của mô hình sẽ tập trung ở các nhận xét 3 sao;
nhận xét ngắn và vừa phải, rất hợp với biểu diễn túi từ.}

\subsection{Biểu diễn văn bản}
Văn bản được chuẩn hóa (chữ thường, bỏ URL, chữ số, ký tự đặc biệt) rồi biểu diễn bằng TF-IDF với unigram và bigram
(Đoạn mã~\ref{code:ch2_text__normalize},~\ref{code:ch2_text__tfidf}). Bigram quan trọng với cảm xúc: ``not good'',
``too small'', ``runs large'' mang nghĩa khác hẳn từng từ riêng lẻ. Từ điển giữ \SL{ch2_comments/models/logreg/vocab} đặc
trưng có tần suất tài liệu ít nhất 3; TF dùng dạng $1+\log(\text{tf})$ để giảm ảnh hưởng của từ lặp nhiều lần. Ma trận
thu được thưa: mỗi nhận xét chỉ có vài chục phần tử khác 0 trên hàng chục nghìn cột, nên mô hình tuyến tính huấn luyện
trong vài giây.

\maNguon{ch2_text__normalize}{Chuẩn hóa văn bản trước khi tách từ}
\maNguon{ch2_text__tfidf}{Bộ biến đổi TF-IDF unigram + bigram}
\hinhpdf{ch2_cmt_wordcloud}{1}{Đám mây từ của nhận xét được giới thiệu (trái) và không được giới thiệu (phải).}{ch2-cmt-cloud}

\subsection{Xây dựng mô hình}
Năm mô hình được so sánh trên cùng biểu diễn TF-IDF (Đoạn mã~\ref{code:ch2_text__text_models}): Logistic Regression có
cân bằng lớp, Multinomial Naive Bayes, Complement Naive Bayes, Linear SVM~\cite{ch2:joachims1998} (được hiệu chuẩn xác suất
bằng \texttt{CalibratedClassifierCV}) và Softmax Regression; baseline luôn dự đoán ``giới thiệu''.

\maNguon{ch2_text__text_models}{Năm mô hình phân loại văn bản trên TF-IDF}
\maNguon{ch2_text__top_terms}{Trích các từ/cụm từ có hệ số lớn nhất của mỗi lớp}

\subsection{Đánh giá mô hình}
\bang{ch2_cmt_val}{Kết quả phân loại nhận xét trên tập Validation}{ch2-cmt-val}
\hinhpdf{ch2_cmt_compare}{1}{Accuracy, macro-F1, F1 lớp ``không giới thiệu'' và ROC-AUC của các mô hình trên Validation.}{ch2-cmt-compare}

\bonhanxet{Logistic Regression (macro-F1 \SL{ch2_comments/val/logreg/macro_f1}), Linear SVM
(\SL{ch2_comments/val/svm/macro_f1}) và Complement NB (\SL{ch2_comments/val/cnb/macro_f1}) gần như ngang nhau, còn
Multinomial NB tụt hẳn (\SL{ch2_comments/val/nb/macro_f1}) vì Recall lớp âm chỉ \SL{ch2_comments/val/nb/rec_neg}.
Complement NB có Recall lớp âm cao nhất (\SL{ch2_comments/val/cnb/rec_neg}).}
{Multinomial NB học từ tần suất từ trong mỗi lớp; lớp âm có ít văn bản nên ước lượng kém và bị lớp dương lấn át. Complement
NB ước lượng trên \emph{phần bù} của lớp nên khắc phục đúng điểm này~\cite{ch2:rennie2003}. Logistic Regression có
\texttt{class\_weight="balanced"} cân bằng Precision và Recall của lớp âm tốt nhất.}
{Với văn bản ngắn và TF-IDF, các mô hình tuyến tính là lựa chọn mạnh, rẻ và dễ diễn giải. Khi lớp quan trọng là lớp thiểu
số, cân bằng lớp quan trọng hơn việc chọn họ mô hình.}

\hinhpdf{ch2_cmt_confusion}{1}{Trái: ma trận nhầm lẫn của Logistic Regression trên tập Test. Phải: Softmax Regression dự đoán số sao
(tỷ lệ theo hàng).}{ch2-cmt-cm}
\hinhpdf{ch2_cmt_terms}{1}{Các từ/cụm từ có hệ số dương (giới thiệu) và âm (không giới thiệu) lớn nhất của Logistic Regression.}{ch2-cmt-terms}

Hình~\ref{fig:ch2-cmt-terms} cho thấy mô hình học được ngôn ngữ cảm xúc có nghĩa: phía âm là các từ thất vọng, trả hàng
và chê độ vừa (disappointed, unflattering, cheap, returned, huge); phía dương là các từ khen ngợi (love, perfect,
great, comfortable, beautiful). Nhiều cụm từ hai từ xuất hiện
trong danh sách, xác nhận lợi ích của bigram.

\subsection{Mở rộng: dự đoán số sao bằng Softmax Regression}
Dự đoán số sao là bài toán 5 lớp có thứ tự, khó hơn nhiều. Softmax Regression trên cùng TF-IDF đạt accuracy
\SL{ch2_comments/rating/acc} và macro-F1 \SL{ch2_comments/rating/macro_f1}, nhưng tỷ lệ dự đoán lệch không quá một sao
là \SL{ch2_comments/rating/within1Pct}\% (Bảng~\ref{tab:ch2-cmt-rating}). Lớp 5 sao dễ nhất (F1
\SL{ch2_comments/rating/f1_per_class/4}), các lớp 1--4 sao khó (F1 từ \SL{ch2_comments/rating/f1_per_class/0} tới
\SL{ch2_comments/rating/f1_per_class/2}), vì ranh giới giữa ``3 sao'' và ``4 sao'' trong lời văn là rất mờ.

\bang{ch2_cmt_rating}{F1 từng lớp của Softmax Regression khi dự đoán số sao trên tập Test}{ch2-cmt-rating}

\subsection{Lựa chọn mô hình và kết quả trên tập Test}
\bang{ch2_cmt_test}{Kết quả phân loại nhận xét trên tập Test}{ch2-cmt-test}

\begin{ketqua}{Kết quả cuối cùng -- Thực nghiệm 3}
Mô hình: \SL{ch2_comments/best/label}\\
Macro-F1 = \SL{ch2_comments/best/macro_f1} · Accuracy = \SL{ch2_comments/best/acc} · ROC-AUC = \SL{ch2_comments/best/auc}\\
Lớp ``không giới thiệu'': F1 = \SL{ch2_comments/best/f1_neg} · Precision = \SL{ch2_comments/best/prec_neg} · Recall = \SL{ch2_comments/best/rec_neg}\\
Baseline: Accuracy = \SL{ch2_comments/test/baseline/acc} · Macro-F1 = \SL{ch2_comments/test/baseline/macro_f1}
\end{ketqua}

Trên Test, Logistic Regression giữ vị trí đầu với macro-F1 \SL{ch2_comments/test/logreg/macro_f1}, hơn Complement NB
(\SL{ch2_comments/test/cnb/macro_f1}) và Linear SVM (\SL{ch2_comments/test/svm/macro_f1}). So với baseline có accuracy
\SL{ch2_comments/test/baseline/accPct}\% nhưng không phát hiện được nhận xét tiêu cực nào, mô hình tìm ra
\SL{ch2_comments/best/rec_negPct}\% số khách hàng không hài lòng. Phần lớn lỗi còn lại là nhận xét 3 sao và nhận xét
mỉa mai, giới hạn cố hữu của biểu diễn túi từ vốn bỏ qua thứ tự từ; Chương~4 sẽ dùng mạng hồi quy để khắc phục.

\subsection{Triển khai hệ thống}
Dịch vụ \texttt{ch2\_comments} nhận một trường văn bản tự do, chạy cùng pipeline chuẩn hóa -- TF-IDF -- mô hình đã huấn
luyện, trả về nhãn ``Giới thiệu''/``Không giới thiệu'' và xác suất của hai lớp. Ví dụ, nhận xét ``The fabric is cheap and the
dress runs very small, I returned it'' được phân loại ``Không giới thiệu'' với xác suất cao.

\maNguon{serve_ch2__predict_comments}{Phân loại một nhận xét khách hàng}

\anhUI{ch2_comments_web}{Giao diện Web App -- phân loại nhận xét khách hàng.}{ch2-cmt-web}
\anhMobile{ch2_comments_mobile}{Giao diện phân loại nhận xét trên điện thoại.}{ch2-cmt-mobile}

% =====================================================================
\section{Tổng kết chương}
\label{sec:ml-summary}
Chương này đi từ nền tảng lý thuyết tới ba hệ thống dự đoán hoàn chỉnh. Ba bài học xuyên suốt:

\begin{enumerate}
  \item \textbf{Chỉ số đánh giá phải khớp với bài toán.} Trên dữ liệu mất cân bằng, baseline đạt accuracy
        \SL{ch2_diabetes/test/baseline/accPct}\% (tiểu đường) và \SL{ch2_comments/test/baseline/accPct}\% (nhận xét) mà
        không có giá trị sử dụng; F1, PR-AUC và macro-F1 mới phản ánh đúng chất lượng.
  \item \textbf{Không có mô hình tốt nhất cho mọi dữ liệu.} Gradient boosting thắng trên dữ liệu y tế có quan hệ ngưỡng
        (F1 \SL{ch2_diabetes/best/f1}); Lasso thắng trên dữ liệu giá nhà gần tuyến tính ($R^2$ \SL{ch2_house/best/r2});
        Logistic Regression trên TF-IDF thắng với văn bản (macro-F1 \SL{ch2_comments/best/macro_f1}). Thiên kiến quy nạp
        của mô hình phải phù hợp với cấu trúc của dữ liệu.
  \item \textbf{Quy trình quan trọng không kém thuật toán.} Làm sạch có hiểu biết miền (giữ ngoại lai có thật), Pipeline
        tránh rò rỉ, chọn ngưỡng trên Validation, chỉ dùng Test một lần, và đóng gói tiền xử lý cùng mô hình khi triển khai.
\end{enumerate}

\begin{table}[H]\centering
\caption{Tóm tắt ba thực nghiệm của Chương 2}\label{tab:ch2-summary}
\small
\begin{tabular}{p{3.2cm}p{3.4cm}p{3.6cm}p{4.2cm}}
\toprule
Thực nghiệm & Kiểu bài toán & Mô hình tốt nhất & Kết quả Test \\
\midrule
Tiểu đường & Phân loại nhị phân, mất cân bằng & \SL{ch2_diabetes/best/label} & F1 \SL{ch2_diabetes/best/f1}, ROC-AUC \SL{ch2_diabetes/best/auc} \\
Giá nhà & Hồi quy & \SL{ch2_house/best/label} & $R^2$ \SL{ch2_house/best/r2}, MAPE \SL{ch2_house/best/mapePct}\% \\
Nhận xét & Phân loại văn bản & Logistic Regression (TF-IDF) & Macro-F1 \SL{ch2_comments/best/macro_f1}, Acc \SL{ch2_comments/best/acc} \\
\bottomrule
\end{tabular}
\end{table}

Giới hạn chung của các mô hình trong chương là chúng làm việc trên \emph{đặc trưng do con người thiết kế}: các cột của bảng
hay túi từ TF-IDF. Với ảnh và chuỗi, việc thiết kế đặc trưng bằng tay trở nên bất khả thi. Hai chương tiếp theo trình bày
cách mạng nơ-ron sâu tự học đặc trưng: mạng tích chập cho ảnh (Chương~3) và mạng hồi quy cho chuỗi (Chương~4).
